% संपूर्ण मराठी ग्रंथातील प्रकरण 69: कट-निर्मूलन
% हा संपादनयोग्य स्रोतखंड आहे. संकलनासाठी संपूर्ण स्रोतसंग्रहातील
% अक्षररूपे, पूर्वभाग, चिन्हव्याख्या आणि आकृती-स्रोत आवश्यक आहेत.
% मूळ: Open Logic Project; मजकूर आणि रूपांतर CC BY 4.0.
% यंत्रानुवाद, दुरुस्ती आणि तपासणी: OpenAI Codex —
% GPT-5.6 Sol आणि GPT-6 Sol, दोन्ही Ultra effort.
% दुरुस्ती-पुनरावलोकन: GPT-6.1 Sol, Ultra effort.
\chapter{कट-निर्मूलन}\label{pt:cut:chap}
\begin{editorial}हे पूर्ण अनुवादित पूरक प्रकरण स्थानिक निदानासाठी मांडले आहे. मूळ माएहारा-सिद्धतेतील गाळलेला चर-पदाचा प्रसंग SOL6-A296 मध्ये भाषा-सीमेच्या आधारे पूर्ण केला आहे; OLINC-344 ही ऐतिहासिक स्रोत-नोंद आहे.\end{editorial}










\section{परिचय}\label{pt:cut:int:sec}

\Cut{} नियम असा आहे:
\[
\Axiom$\Gamma \fCenter \Delta, A$
\Axiom$A, \Pi \fCenter \Lambda$
\RightLabel{\Cut}
\BinaryInf$\Gamma, \Pi \fCenter \Delta, \Lambda$
\DisplayProof
\]
(सूत्र~$A$ याला
\emph{कट-सूत्र} म्हणतात.)

\Cut{} हा नियम गेंटझेन यांच्या
मूळ क्रमवर्ती कलनात~\Log{LK} होता.
गेंटझेन यांचे \emph{Hauptsatz}
(“मुख्य प्रमेय”) असे सांगते की
~\Log{LK} मध्ये सिद्धतायोग्य असलेली
प्रत्येक क्रमवर्ती \Cut{} न वापरताही
सिद्धतायोग्य आहे; म्हणजे
~\Log{LK} मधील सिद्धतांतून
\Cut{} “काढून टाकता” येतो.
आपल्या \Log{G1c} आणि \Log{G3c}
पद्धतींत \Cut{} नाही; तरी त्यांच्याबाबतही
प्रश्न पडतो: \Cut{} आणि
~\Log{G1c} चे (किंवा \Log{G3c} चे)
इतर नियम वापरणाऱ्या सिद्धता मधून
\Cut{} काढून टाकता येईल का?
आपल्या आधी निश्चित केलेल्या
परिभाषेत हाच प्रश्न असा आहे:
\Cut{} हा \Log{G1c} चा
(किंवा \Log{G3c} चा)
ग्राह्य नियम आहे का?

कट-निर्मूलन प्रमेय हा सिद्धता
उपपत्तीतील कदाचित सर्वाधिक महत्त्वाचा
आणि प्रसिद्ध निष्कर्ष आहे.
सुरुवातीला आवश्यक वाटू शकणारा
नियम \Log{G1c} आणि \Log{G3c}
यांना प्रत्यक्षात लागत नाही, हे
तो दाखवतो. कारण प्रत्यक्ष सिद्धता
औपचारिक करताना उपप्रमेयांचा
वापर कसा दाखवायचा हा प्रश्न
पडतो. गणितज्ञ प्रथम एक निष्कर्ष
सिद्ध करतात आणि मग दुसऱ्या
निष्कर्षाच्या सिद्धतेत त्याचा
आधार घेतात. म्हणून उपप्रमेय
$L$ ची सिद्धता प्रथम
$\Sequent L$ या क्रमवर्तीच्या
सिद्धता म्हणून औपचारिक करता येईल.
मग $L$ वापरणाऱ्या दुसऱ्या
निष्कर्ष~$R$ ची सिद्धता
$L \Sequent
R$ या क्रमवर्तीच्या
सिद्धता म्हणून औपचारिक करू.
केवळ $R$ ची, म्हणजे
$\Sequent R$ ची
सिद्धता आहे, हे कसे ठरवायचे?
त्या दोन सिद्धता जोडण्याचा
मार्ग हवा; \Cut{} नियमाने
ते करता येईल.

\Log{G1c} आणि \Log{G3c}
संपूर्ण आहेत, असे आपण सांगितले:
म्हणून सिद्ध होण्याची अपेक्षा
असलेल्या सर्व वैध क्रमवर्ती
ते सिद्ध करतात. हे थेटही
सिद्ध करता येते. पण गेंटझेन
लिहीत होते तेव्हा केवळ स्वयंसिद्धकीय
निष्पत्ती-पद्धती संपूर्ण आहेत,
हे ज्ञात होते. \Log{LK} ची
संपूर्णता दाखवण्यासाठी गेंटझेन
यांनी स्वयंसिद्धकीय पद्धतीतील
निष्पत्तीचे
~\Log{LK} मधील सिद्धता मध्ये
भाषांतर दिले; त्यात
\Cut{} अत्यावश्यक होता.
कट-निर्मूलन केवळ अभिजात
तर्कशास्त्रापुरते महत्त्वाचे नाही:
इतर अनेक तर्कशास्त्रांनाही
क्रमवर्ती कलने आहेत. काहींसाठी
त्या कलनांची संपूर्णता थेट
सिद्ध करणे कठीण किंवा अशक्य
असते; त्यामुळे इतर पद्धतींमधून
\Cut{} वापरून केलेले भाषांतर
हाच संपूर्णता दाखवण्याचा मार्ग असतो.
मग \Cut{} अनावश्यक आहे आणि
\Cut{} नसलेली पद्धतही संपूर्ण
आहे, हे स्थापित करण्यासाठी
सिद्धता-उपपत्तीय कट-निर्मूलन
सिद्ध करणे आवश्यक ठरते.

कट-निर्मूलनाचा उपयोग इतर
स्वतंत्रपणे महत्त्वाचे निष्कर्ष
मिळवण्यासाठीही होतो. त्यांपैकी
क्रेगचे अंतर्वेशन उपप्रमेय आणि
एरब्राँचे प्रमेय ही दोन उदाहरणे
आपण पाहू.

कट-निर्मूलनाच्या सिद्धतेत
बहुधा \Cut{} वापरणाऱ्या
सिद्धता चे त्याविना
सिद्धतेत रूपांतर कसे करावे,
हे दाखवले जाते. ही रूपांतरे
सहसा “स्थानिक” असतात:
सिद्धतेचा एक भाग
(साधारणपणे \Cut{} अनुमानाने
संपणारी उप-सिद्धता)
घेऊन त्याच्या जागी
त्याच क्रमवर्तीची दुसरी
उप-सिद्धता ठेवतो;
तिच्यातील \Cut{} अनुमान
काढलेले किंवा इतर
अनुमानांनी बदललेले असते.
अशा युक्तिवादांनी \Cut{}
असलेल्या सिद्धता चे
त्याविना सिद्धतांत रूपांतर
करणारी पायरी-पायरीची
अल्गोरिदम मिळते. उपयोगांत
कट-निर्मूलन प्रक्रिया वापरताना
यातून अधिक माहिती मिळते.
उदाहरणार्थ, अंतर्वेशन उपप्रमेयाची
अशी प्रतिमान-उपपत्तीय सिद्धता
आहे की “$A \lif B$”
वैध असेल, तर अंतर्वेशक अस्तित्वात
असतो (आत्ता अंतर्वेशक म्हणजे
काय याकडे दुर्लक्ष करा).
कट-निर्मूलन वापरणारा
सिद्धता-उपपत्तीय युक्तिवाद
$A
\lif B$ ची सिद्धता
घेऊन अंतर्वेशक देणारी
अल्गोरिदम पुरवतो. प्रतिमान
उपपत्तीच्या युक्तिवादाच्या
तुलनेत ही सिद्धता-उपपत्तीय
पद्धत \emph{रचनात्मक} आहे.

कट-निर्मूलन प्रमेय सिद्ध
करण्याच्या दोन प्रचलित पद्धती
आहेत. पहिली, गेंटझेनपासून
आलेली, सिद्धतेतील सर्वोच्च
स्थानी असलेले कट काढता
येतात हे दाखवते आणि मग
असे कट संपेपर्यंत काढत जाते.
दुसरी, मूळची श्यूटे आणि
टेट यांची, सिद्धता मधील
सर्वाधिक गुंतागुंतीच्या कटांवर
लक्ष केंद्रित करते; इतर
कोणत्याही कटच्या सूत्र
ची खोली अधिक नाही, या
अर्थाने महत्तम असलेले
कट ती एकामागोमाग काढते.
दोन्ही पद्धतींचे फायदे
आणि मर्यादा आहेत. काही
पद्धतींमध्ये एक मार्ग
चालतो, तर दुसरा अवघड
किंवा अशक्य असतो.
म्हणून आपण दोन्ही मार्ग
वर्णन करू. ~\Log{G3c}
साठी कट-निर्मूलन सिद्ध करू.
खरे तर नेहमीचा \Cut{}
नव्हे, तर \emph{संदर्भ-सामायिक कट}~\CutCS:
\[
\Axiom$\Gamma \fCenter \Delta, A$
\Axiom$A, \Gamma \fCenter \Delta$
\RightLabel{\CutCS}
\BinaryInf$\Gamma \fCenter \Delta$
\DisplayProof
\]
एखाद्या क्रमवर्ती कलनात
दुर्बलीकरण आणि आकुंचन
हे प्रत्यक्ष नियम म्हणून
(जसे \Log{G1c} मध्ये)
किंवा ग्राह्य नियम म्हणून
(जसे \Log{G3c} मध्ये)
असतील, तर \Cut{} द्वारे
\CutCS चे अनुकरण करता
येते, आणि उलटही.
आपण \Cut ऐवजी \CutCS{}
निवडतो, कारण पहिली
कट-निर्मूलन पद्धत \CutCS साठी
वर्णन करणे सोपे आहे,
आणि दुसरी पद्धत केवळ
\CutCS साठीच चालते.













\section{सर्वोच्च स्थानी असलेले कट काढणे}\label{pt:cut:top:sec}

\begin{defn}
एका \CutCS{} नियमाने संपणारी सिद्धता~$\pi$
विचारात घ्या:
\[
\AxiomC{}
\RightLabel{$\pi_1$}
\Deduce$\Gamma \fCenter \Delta, A$
\AxiomC{}
\RightLabel{$\pi_2$}
\Deduce$A, \Gamma \fCenter \Delta$
\RightLabel{\CutCS}
\BinaryInf$\Gamma \fCenter \Delta$
\DisplayProof
\]
तिच्यात अन्य \CutCS{} अनुमाने नसावीत.
$\pi$ ची \emph{कट-उंची} म्हणजे
$\cheight{\pi} = \pheight{\pi_1} + \pheight{\pi_2}$.
$\pi$ ची \emph{कट-कोटी} म्हणजे
$\cutrank{\pi} = \depth{A}$.
\end{defn}

\begin{lem}\label{pt:cut:top:lem:cut-adm-G3c}
\CutCS{} नियम~\Log{G3c} मध्ये ग्राह्य आहे.
दुसऱ्या शब्दांत, सिद्धता~$\pi$ ही एका
\CutCS{} ने संपत असेल आणि इतरत्र
फक्त~$\Log{G3c}$ चे नियम वापरत असेल,
तर त्याच अंतिम क्रमवर्तीची
\Log{G3c} मधील सिद्धता~$\pi'$ असते.
\end{lem}

\begin{proof}
कटची उंची आणि कोटी यांवरील
दुहेरी विगमनाने सिद्धता करू.
सिद्धता~$\pi$ चा शेवट असा आहे:
\[
\AxiomC{}
\RightLabel{$\pi_1$}
\Deduce$\Gamma \fCenter \Delta, A$
\AxiomC{}
\RightLabel{$\pi_2$}
\Deduce$A, \Gamma \fCenter \Delta$
\RightLabel{\CutCS}
\BinaryInf$\Gamma \fCenter \Delta$
\DisplayProof
\]

विगमनाचा आधार: $\cheight{\pi} = 0$
आणि $\cutrank{\pi} = 0$.
$\cheight{\pi} = \pheight{\pi_1}
+ \pheight{\pi_2} = 0$ असल्याने
$\Gamma \Sequent \Delta, A$ आणि
$A, \Gamma \Sequent \Delta$
ही दोन्ही स्वयंसिद्धके आहेत.
दोन प्रसंग आहेत: (a)~$A$ हा
त्यांपैकी एखाद्यात मुख्य नाही,
किंवा (b)~$A$ हा दोन्हींत मुख्य आहे.
खाली दोन्ही प्रसंगांत
$\Gamma \Sequent \Delta$ हे स्वयंसिद्धक
असलेच पाहिजे हे दाखवू
(त्यासाठी विगमन गृहीतक लागत नाही).

आता $\pi_1$ आणि $\pi_2$ स्वयंसिद्धके
आहेत की अनुमानाने संपतात, तसेच त्या
अनुमानात~$A$ मुख्य आहे की नाही,
यावरून प्रसंग वेगळे करू. प्रसंग A व~B
विगमनाचा आधार समाविष्ट करतात.
विगमन पायरीत $\cheight{\pi} > 0$
किंवा $\cutrank{\pi} > 0$
(किंवा दोन्ही) असे मानतो. विगमन
गृहीतकानुसार, एका \CutCS{} ने
संपणाऱ्या सिद्धता मधील अंतिम
\CutCS{} काढता येतो, जर तिची
कट-कोटी $= \cutrank{\pi}$ आणि
कटची उंची $< \cheight{\pi}$
असेल, किंवा तिची कट-कोटी
$< \cutrank{\pi}$ असेल.
$\cheight{\pi} = 0$ चा प्रसंग
(दोन्ही आधार स्वयंसिद्धके) A व~B
मध्ये येतो. $\cheight{\pi} > 0$
चे प्रसंग C--D मध्ये येतात.

\paragraph{प्रसंग A:}
काही $B$ हे $\Gamma$ आणि~$\Delta$
या दोन्हींत येते. $B$ अणुसूत्र
असेल, तर $\Gamma \Sequent \Delta$
हे स्वयंसिद्धक असते; अन्यथा
\ref{pt:seq:ex:prop:G3c-axioms} नुसार
त्याची \Log{G3c} मधील \CutCS{}
विरहित सिद्धता~$\pi'$ असते.
यात विगमन आधाराचा प्रसंग~(a) येतो:
$\Gamma \Sequent \Delta, A$
किंवा $A, \Gamma \Sequent \Delta$
हे स्वयंसिद्धक असून त्यात $A$
मुख्य नसेल, तर अन्य अणुसूत्र~$B$
हे $\Gamma$ आणि~$\Delta$ या
दोन्हींत असते. विगमन पायरीतही
एखादा आधार स्वयंसिद्धक असून त्यात
$A$ मुख्य नसेल, तर हा प्रसंग
लागू होतो; दुसरा आधार नियमाचा
निष्कर्ष असला तरी चालते.

\paragraph{प्रसंग B:}
दोन्ही आधार स्वयंसिद्धके आहेत
आणि $A$ मुख्य, म्हणून अणुसूत्र,
आहे. डावा आधार
$\Gamma \Sequent \Delta, A$
पाहा. या $\Gamma \Sequent \Delta, A$
मध्ये $A$ मुख्य असल्यामुळे ते
$\Gamma$ मध्ये असलेच पाहिजे
(नसते तर ही क्रमवर्ती स्वयंसिद्धक
ठरली नसती). त्याचप्रमाणे,
$A$ हे $\Delta$ मध्येही असलेच
पाहिजे; अन्यथा
$A, \Gamma \Sequent \Delta$
स्वयंसिद्धक ठरणार नाही.
त्यामुळे $A$ अणुसूत्र असून
$\Gamma$ आणि $\Delta$ या
दोन्हींत येते; म्हणजे
$\Gamma \Sequent \Delta$
हे स्वयंसिद्धक आहे. हा
विगमन आधाराचा प्रसंग~(b) आहे.

\begin{center}
प्रसंग A आणि B चे पर्यायी रूप
\end{center}

\paragraph{प्रसंग A:}
कटच्या आधारांपैकी एक
$C, \Gamma' \Sequent \Delta', C$
या रूपाचे स्वयंसिद्धक आहे,
जेथे $C$ अणुसूत्र आहे.
$C$ हे कटचे सूत्र~$A$
नसेल, तर $C$ हे $\Gamma$
आणि~$\Delta$ या दोन्हींत
असलेच पाहिजे. मग
$\Gamma \Sequent \Delta$
हे स्वयंसिद्धक असून तेच~$\pi'$
घेता येते. $C \ident A$
कटचे सूत्र असेल, तर
डावा आधार स्वयंसिद्धक असल्यास
$A \in \Gamma$ आणि
$\Gamma = \Gamma', A$;
उजवा आधार स्वयंसिद्धक असल्यास
$A \in \Delta$ आणि
$\Delta = \Delta', A$.
पहिल्या प्रसंगात $\pi_2$
ही $A, A, \Gamma' \Sequent \Delta$
ची सिद्धता असते;
\LeftR{\Contraction} ची ग्राह्यता
(\ref{pt:seq:inv:prop:G3c-cont-adm})
वापरून $\Gamma \Sequent \Delta$
ची सिद्धता मिळते.
दुसऱ्या प्रसंगात $\pi_1$
ही $\Gamma \Sequent \Delta', A, A$
ची सिद्धता असते;
\RightR{\Contraction} ची ग्राह्यता
वापरून $\Gamma \Sequent \Delta$
ची सिद्धता मिळते.


\paragraph{प्रसंग B:}
एखादा आधार
$\lfalse, \Gamma' \Sequent \Delta'$
या रूपाचा स्वयंसिद्धक आहे.
डावा आधार असा असेल, तर
$\lfalse \in \Gamma$ आणि
$\Gamma \Sequent \Delta$
हेही स्वयंसिद्धक आहे. उजवा
आधार असा असेल, तर
$\lfalse \in \Gamma$
(मग पुन्हा $\Gamma \Sequent \Delta$
स्वयंसिद्धक असते) किंवा
$A \ident \lfalse$.
दुसऱ्या उपप्रसंगात $\pi_1$
ही $\Gamma \Sequent \Delta, \lfalse$
ची सिद्धता आहे. $\pi_1$
मधील प्रत्येक क्रमवर्तीच्या
उत्तरांगातून $\lfalse$ ची
एक प्रत काढून
$\Gamma \Sequent \Delta$
ची सिद्धता मिळते.
(सराव: $\pheight{\pi_1}$
वरील विगमनाने हे तपासा.)

आता प्रसंग~A किंवा B लागू
होत नाही असे समजा. उर्वरित
प्रसंगांत $\cheight{\pi} > 0$;
म्हणजे किमान एका आधाराचा
शेवट नियमाच्या अनुमानाने होतो.

\paragraph{प्रसंग C आणि D: कटचे स्थान बदलणे.}
प्रसंग~C आणि~D मधील सर्व शक्यतांत
एक सामाईक नमुना आहे. सुरुवातीला
\CutCS{} ने संपणारी सिद्धता असते;
तिच्या एका आधाराचा निष्कर्ष नियम~$R$
ने मिळतो, पण कटचे सूत्र~$A$
त्या नियमात मुख्य नसते (काही अन्य
सूत्र~$D$ मुख्य असते).
या सिद्धता मध्ये \CutCS{}
हा नियम~$R$ नंतर येतो; तिचे
योजनात्मक रूप असे:
\[
\AxiomC{}
\RightLabel{$\pi_1$}
\Deduce$\Gamma \fCenter \Delta', D, A$
\AxiomC{}
\RightLabel{$\pi_3$}
\Deduce$A, \Gamma, \Pi \fCenter \Delta', \Lambda$
\RightLabel{$R$}
\UnaryInf$A, \Gamma \fCenter \Delta', D$
\RightSubproofLabel{$\pi_2$}
\RightLabel{\CutCS}
\BinaryInf$\Gamma \fCenter \Delta', D$
\DisplayProof
\]
हे रूप फक्त त्या प्रसंगाचे आहे
जेथे $R$ हा एक-आधाराचा उजवा
नियम आहे, $A$ हे $R$ चे
मुख्य सूत्र नाही, आणि
\emph{मुख्य} असलेले सूत्र~$D$
हे \CutCS{} च्या उजव्या
आधाराच्या उत्तरांगात आहे.
$D$ पूर्वांगात असेल (म्हणजे
$R$ डावा नियम असेल) किंवा
\CutCS{} च्या डाव्या आधारात
असेल, तर रूप अर्थातच बदलेल.
$\Pi$ आणि~$\Lambda$ मधील
सूत्रे या नियम~$R$
च्या सक्रिय सूत्रे दर्शवतात.

आता हा क्रम उलटून
$R$ हा~\CutCS{} नंतर येणारी
सिद्धता विचारात घ्या:
\[
\AxiomC{}
\RightLabel{$\pi_1'$}
\Deduce$\Gamma, \Pi \fCenter \Delta', D, \Lambda, A$
\AxiomC{}
\RightLabel{$\pi_3'$}
\Deduce$A, \Gamma, \Pi \fCenter \Delta', D, \Lambda$
\RightLabel{\CutCS}
\BinaryInf$\Gamma, \Pi \fCenter \Delta', D, \Lambda$
\RightLabel{$R$}
\UnaryInf$\Gamma, \fCenter \Delta', D, D$
\DisplayProof
\]
$\pi_1$ आणि $\pi_3$ यांवर
उंची न वाढवणारे दुर्बलीकरण
लावून अनुक्रमे $\pi_1'$ आणि
$\pi_3'$ मिळवतो.

\CutCS{} आणि~$R$ यांचा क्रम
बदलत असल्यामुळे याला
“कटला~$R$ च्या वर नेणे” म्हणतात.
याने \CutCS{} च्या उजव्या
आधारावर संपणाऱ्या सिद्धता
ची उंची कमी होते; म्हणून कटची
उंची कमी होते. म्हणजे
नव्या \CutCS{} वरील
उप-सिद्धता ची उंची
$\pi_1$ आणि $\pi_3$ यांच्या
उंचीपेक्षा मोठी नसते, म्हणून
कटची उंची कमी होते
(उजवीकडील एक अनुमान काढले आहे).
आता विगमन गृहीतक लावून \CutCS{}
अनुमान काढता येते. शेवटी
$\RightR{\Contraction}$ ची ग्राह्यता
वापरून~$D$ ची अतिरिक्त
प्रत काढतो.

\paragraph{प्रसंग C:}
$\pi_1$ आणि $\pi_2$ यांपैकी
किमान एकाचा शेवट एक-आधाराच्या
नियम~$R$ ने होतो, आणि त्यात
$A$ मुख्य नसते. एकूण 18 प्रसंग
आहेत: \CutCS{} च्या डाव्या की
उजव्या आधारात $A$ मुख्य नाही
आणि नऊ एक-आधाराच्या नियमांपैकी
(\LeftR{\lnot}, \RightR{\lnot},
\RightR{\lor}, \LeftR{\land},
\RightR{\lif}, आणि चार संख्यक नियम)
$R$ कोणता आहे, यावर ते अवलंबून
असतात. आपण फक्त दोन उदाहरणे
पाहू: $\pi_2$ च्या अंतिम
अनुमानात $A$ मुख्य नाही आणि ते
अनुमान $\RightR{\lif}$ किंवा
\LeftR{\lexists} आहे.
इतर प्रसंग यांसारखेच असून
सरावासाठी ठेवतो.

समजा $\pi$ चा शेवट असा आहे:
\[
\AxiomC{}
\RightLabel{$\pi_1$}
\Deduce$\Gamma \fCenter \Delta', B \lif C, A$
\AxiomC{}
\RightLabel{$\pi_3$}
\Deduce$A, B, \Gamma \fCenter \Delta', C$
\RightLabel{$\RightR{\lif}$}
\UnaryInf$A, \Gamma \fCenter \Delta', B \lif C$
\RightSubproofLabel{$\pi_2$}
\RightLabel{\CutCS}
\BinaryInf$\Gamma \fCenter \Delta', B \lif C$
\DisplayProof
\]
येथे $\Delta = \Delta', B \lif C$.
$\pi_1$ ला उंची न वाढवणारे
दुर्बलीकरण
(\ref{pt:seq:adm:prop:weak-G3c-adm})
लावून $B, \Gamma \fCenter
\Delta', B \lif C, C, A$
ची सिद्धता $\pi_1'$ मिळवतो
(डावीकडे $B$ आणि उजवीकडे
$C$ घालतो). त्याचप्रमाणे
\ref{pt:seq:adm:prop:weak-G3c-adm}
हे~$\pi_3$ ला लावून
$A, B, \Gamma \fCenter
\Delta', B \lif C, C$
ची सिद्धता~$\pi_3'$ मिळते
(उजवीकडे $B \lif C$ घालतो).
विगमन गृहीतक पुढील सिद्धतेला
लागू होते:
\[
\AxiomC{}
\RightLabel{$\pi_1'$}
\Deduce$B, \Gamma \fCenter \Delta', B \lif C, C, A$
\AxiomC{}
\RightLabel{$\pi_3'$}
\Deduce$A, B, \Gamma \fCenter \Delta', B \lif C, C$
\RightLabel{\CutCS}
\BinaryInf$B, \Gamma \fCenter \Delta', B \lif C, C$
\DisplayProof
\]
यातील कटचे सूत्र~$A$
आहे: या सिद्धता ची कट-कोटी
$\pi$ इतकीच, पण कटची उंची
कमी आहे. त्यामुळे \CutCS{}
विरहित \Log{G3c} मधील
$B, \Gamma \fCenter
\Delta', B \lif C, C$
ची सिद्धता~$\pi_4$ मिळते.
$\RightR{\lif}$ लावून
$\Gamma \fCenter \Delta',
B \lif C, B \lif C$
ची सिद्धता मिळते:
\[
\AxiomC{}
\RightLabel{$\pi_4$}
\Deduce$B, \Gamma \fCenter \Delta', B \lif C, C$
\RightLabel{$\RightR{\lif}$}
\UnaryInf$\Gamma \fCenter \Delta', B \lif C, B \lif C$
\DisplayProof
\]
$\Gamma \Sequent \Delta$,
म्हणजे $\Gamma \fCenter
\Delta', B \lif C$,
याची सिद्धता~$\pi'$
मिळवण्यासाठी
\ref{pt:seq:inv:prop:G3c-cont-adm},
म्हणजे आकुंचनाची ग्राह्यता,
वापरा.

आता $\pi$ चा शेवट
पुढीलप्रमाणे आहे असे समजा:
\[
\AxiomC{}
\RightLabel{$\pi_1$}
\Deduce$\lexists[x][B(x)], \Gamma' \fCenter \Delta, A$
\AxiomC{}
\RightLabel{$\pi_3$}
\Deduce$A, B(c), \Gamma' \fCenter \Delta$
\RightLabel{$\LeftR{\lexists}$}
\UnaryInf$A, \lexists[x][B(x)], \Gamma' \fCenter \Delta$
\RightSubproofLabel{$\pi_2$}
\RightLabel{\CutCS}
\BinaryInf$\lexists[x][B(x)], \Gamma' \fCenter \Delta$
\DisplayProof
\]
पुन्हा पुढील सिद्धतेला
विगमन गृहीतक लावू:
\[
\AxiomC{}
\RightLabel{$\pi_1[B(c) \Sequent {}]$}
\Deduce$\lexists[x][B(x)], B(c), \Gamma' \fCenter \Delta, A$
\AxiomC{}
\LeftLabel{$\pi_3[\lexists[x][B(x)] \Sequent {}]$}
\Deduce$A, \lexists[x][B(x)], B(c), \Gamma' \fCenter \Delta$
\RightLabel{\CutCS}
\BinaryInf$\lexists[x][B(x)], B(c), \Gamma' \fCenter \Delta$
\RightLabel{\LeftR{\lexists}}
\UnaryInf$\lexists[x][B(x)], \lexists[x][B(x)], \Gamma' \fCenter \Delta$
\DisplayProof
\]
येथे आयगेन चराची अट पूर्ण होते.
कारण $\pi_2$ मधील
\LeftR{\lexists} वरील
आयगेन चराच्या अटीनुसार
$c$ हे $\lexists[x][B(x)]$,
$\Gamma'$ किंवा~$\Delta$
यांत येत नाही. शेवटी
\ref{pt:seq:inv:prop:G3c-cont-adm}
लावा.

\paragraph{प्रसंग D:}
$\pi_1$ आणि $\pi_2$ यांपैकी
किमान एकाचा शेवट दोन-आधारांच्या
नियम~$R$ ने होतो आणि त्यात
$A$ मुख्य नसते. असे सहा
प्रसंग आहेत. $\pi_2$ चा
शेवट~$\RightR{\land}$ ने होत
असताना त्याच्या अंतिम अनुमानात
$A$ मुख्य नसतो हा प्रसंग
पाहू; इतर तसेच आहेत.

समजा $\pi$ चा शेवट असा आहे:
\[
\AxiomC{}
\RightLabel{$\pi_1$}
\Deduce$\Gamma \fCenter \Delta', B \land C, A$
\AxiomC{}
\RightLabel{$\pi_3$}
\Deduce$A, \Gamma \fCenter \Delta', B$
\AxiomC{}
\RightLabel{$\pi_4$}
\Deduce$A, \Gamma \fCenter \Delta', C$
\RightLabel{$\RightR{\land}$}
\BinaryInf$A, \Gamma \fCenter \Delta', B \land C$
\RightSubproofLabel{$\pi_2$}
\RightLabel{\CutCS}
\BinaryInf$\Gamma \fCenter \Delta$
\DisplayProof
\]
विगमन गृहीतक पुढील
सिद्धता ना लागू होते:
\[
\AxiomC{}
\RightLabel{$\pi_1'$}
\Deduce$\Gamma \fCenter \Delta', B \land C, B, A$
\AxiomC{}
\RightLabel{$\pi_3'$}
\Deduce$A, \Gamma \fCenter \Delta', B \land C, B$
\RightLabel{\CutCS}
\BinaryInf$\Gamma \fCenter \Delta', B \land C, B$
\DisplayProof
\]
आणि
\[
\AxiomC{}
\RightLabel{$\pi_1''$}
\Deduce$\Gamma \fCenter \Delta', B \land C, C, A$
\AxiomC{}
\RightLabel{$\pi_4'$}
\Deduce$A, \Gamma \fCenter \Delta', B \land C, C$
\RightLabel{\CutCS}
\BinaryInf$\Gamma \fCenter \Delta', B \land C, C$
\DisplayProof
\]
येथे \Log{G3c} मधील
सिद्धता $\pi_1'$ आणि
$\pi_1''$ या $\pi_1$
पासून उंची न वाढवणाऱ्या
दुर्बलीकरणाने
(\ref{pt:seq:adm:prop:weak-G3c-adm})
मिळतात: एकदा उजवीकडे $B$
आणि एकदा $C$ घालून.
सिद्धता $\pi_3'$ आणि
$\pi_4'$ या अनुक्रमे $\pi_3$
आणि $\pi_4$ पासून
उंची न वाढवणाऱ्या
दुर्बलीकरणाने मिळतात:
यात $\pi_3$ व $\pi_4$ यांची
उंची टिकवून उजवीकडे $B \land C$
घालून.


विगमन गृहीतकानुसार
\Log{G3c} मधील
$\Gamma \fCenter \Delta', B \land C, B$
आणि
$\Gamma \fCenter \Delta', B \land C, C$
यांच्या सिद्धता~$\pi_5$ आणि~$\pi_6$
मिळतात. $\RightR{\land}$
नियमाने त्यांना जोडून
$\Gamma \Sequent \Delta',
B \land C, B \land C$
ची सिद्धता मिळते:
\[
\AxiomC{}
\RightLabel{$\pi_5$}
\Deduce$\Gamma \fCenter \Delta', B \land C, B$
\AxiomC{}
\RightLabel{$\pi_6$}
\Deduce$\Gamma \fCenter \Delta', B \land C, C$
\RightLabel{\RightR{\land}}
\BinaryInf$\Gamma \fCenter \Delta', B \land C, B \land C$
\DisplayProof
\]
$\Gamma \Sequent \Delta$,
म्हणजे $\Gamma \fCenter
\Delta', B \land C$,
याची सिद्धता मिळवण्यासाठी
\ref{pt:seq:inv:prop:G3c-cont-adm},
म्हणजे आकुंचनाची ग्राह्यता,
लावा.


\paragraph{प्रसंग E: कटची कोटी कमी करणे.}
शेवटचा प्रसंग असा की $\pi_1$
आणि $\pi_2$ दोन्हींचा शेवट
अशा अनुमानांनी होतो ज्यांत
$A$ मुख्य असते. $A$ च्या
संभाव्य मुख्य संकारक नुसार
सहा प्रसंग होतात.

प्रत्येक प्रसंगात कटचे
सूत्र~$A$ असलेले
\CutCS{} अनुमान आपण $A$
च्या उप-सूत्रे वरील
एक किंवा अधिक कटांत बदलतो;
म्हणजे ज्या अनुमानांत $A$
हे मुख्य सूत्र असते,
त्यांच्या सक्रिय सूत्रे
वर कट करतो. म्हणून या
प्रसंगांना कट-अनुमानाचे
“न्यूनीकरण” किंवा “परिवर्तन”
म्हणतात.

\begin{enumerate}
\item $A \ident \lnot B$
असेल, तर $\pi$ चा शेवट असा:
\[
\AxiomC{}
\RightLabel{$\pi_1'$}
\Deduce$B, \Gamma \fCenter \Delta$
\RightLabel{\RightR{\lnot}}
\UnaryInf$\Gamma \fCenter \Delta, \lnot B$
\LeftSubproofLabel{$\pi_1$}
\AxiomC{}
\RightLabel{$\pi_2'$}
\Deduce$\Gamma \fCenter \Delta, B$
\RightLabel{\LeftR{\lnot}}
\UnaryInf$\lnot B, \Gamma \fCenter \Delta$
\RightSubproofLabel{$\pi_2$}
\RightLabel{\CutCS}
\BinaryInf$\Gamma \fCenter \Delta$
\DisplayProof
\]
विगमन गृहीतकानुसार पुढील
सिद्धतेतून \CutCS{} काढता येतो:
\[
\AxiomC{}
\RightLabel{$\pi_2'$}
\Deduce$\Gamma \fCenter \Delta, B$
\AxiomC{}
\RightLabel{$\pi_1'$}
\Deduce$B, \Gamma \fCenter \Delta$
\RightLabel{\CutCS}
\BinaryInf$\Gamma \fCenter \Delta$
\DisplayProof
\]

\item
$A \ident (B \lor C)$ असेल,
तर $\pi$ चा शेवट असा आहे:
\[
\AxiomC{}
\RightLabel{$\pi_1'$}
\Deduce$\Gamma \fCenter \Delta, B, C$
\RightLabel{\RightR{\lor}}
\UnaryInf$\Gamma \fCenter \Delta, B \lor C$
\LeftSubproofLabel{$\pi_1$}
\AxiomC{}
\RightLabel{$\pi_2'$}
\Deduce$B, \Gamma \fCenter \Delta$
\AxiomC{}
\RightLabel{$\pi_2''$}
\Deduce$C, \Gamma \fCenter \Delta$
\RightLabel{\LeftR{\lor}}
\BinaryInf$B \lor C, \Gamma \fCenter \Delta$
\RightSubproofLabel{$\pi_2$}
\RightLabel{\CutCS}
\BinaryInf$\Gamma \fCenter \Delta$
\DisplayProof
\]
\CutCS{} ने संपणारी पुढील
सिद्धता विचारात घ्या:
\[
\AxiomC{}
\RightLabel{$\pi_1'$}
\Deduce$\Gamma \fCenter \Delta, B, C$
\AxiomC{}
\RightLabel{$\pi_2'''$}
\Deduce$C, \Gamma \fCenter \Delta, B$
\RightLabel{\CutCS}
\BinaryInf$\Gamma \fCenter \Delta, B$
\DisplayProof
\]
येथे $\pi_2'''$ ही $\pi_2''$
पासून उंची न वाढवणारे
दुर्बलीकरण लावून मिळते.
$\depth{C} < \depth{B \lor C}$
असल्यामुळे तिची कट-कोटी
$< \cutr{\pi}$ आहे.
विगमन गृहीतकानुसार
$\Gamma \fCenter \Delta, B$
ची \Log{G3c} मधील
सिद्धता~$\pi_1''$ मिळते.
आता \CutCS{} ने संपणारी
पुढील सिद्धता पाहा:
\[
\AxiomC{}
\RightLabel{$\pi_1''$}
\Deduce$\Gamma \fCenter \Delta, B$
\AxiomC{}
\RightLabel{$\pi_2'$}
\Deduce$B, \Gamma \fCenter \Delta$
\RightLabel{\CutCS}
\BinaryInf$\Gamma \fCenter \Delta$
\DisplayProof
\]
हिची कट-कोटी
$= \depth{B} < \depth{B \lor C}$
असल्यामुळे विगमन गृहीतक
पुन्हा लागू होते. त्यातून
$\Gamma \fCenter \Delta$
ची सिद्धता $\pi'$ मिळते.

\item $A \ident (B \land C)$.
हा प्रसंग सरावासाठी आहे.

\item $A \ident (B \lif C)$
असेल, तर $\pi$ चा शेवट असा:
\[
\AxiomC{}
\RightLabel{$\pi_1'$}
\Deduce$B, \Gamma \fCenter \Delta, C$
\RightLabel{\RightR{\lif}}
\UnaryInf$\Gamma \fCenter \Delta, B \lif C$
\LeftSubproofLabel{$\pi_1$}
\AxiomC{}
\RightLabel{$\pi_2'$}
\Deduce$\Gamma \fCenter \Delta, B$
\AxiomC{}
\RightLabel{$\pi_2''$}
\Deduce$C, \Gamma \fCenter \Delta$
\RightLabel{\LeftR{\lif}}
\BinaryInf$B \lif C, \Gamma \fCenter \Delta$
\RightSubproofLabel{$\pi_2$}
\RightLabel{\CutCS}
\BinaryInf$\Gamma \fCenter \Delta$
\DisplayProof
\]
$\pi_2''$ च्या पूर्वांगात $B$ घालून मिळालेल्या
सिद्धतेला $\pi_2'''$ लिहू. आता दोन्ही आधारांचा सामायिक
संदर्भ $B,\Gamma$ आहे. \CutCS{} ने संपणाऱ्या
पुढील सिद्धता ची कट-कोटी $\pi$ पेक्षा कमी आहे:
\[
\AxiomC{}
\RightLabel{$\pi_1'$}
\Deduce$B, \Gamma \fCenter \Delta, C$
\AxiomC{}
\RightLabel{$\pi_2'''$}
\Deduce$C, B, \Gamma \fCenter \Delta$
\RightLabel{\CutCS}
\BinaryInf$B, \Gamma \fCenter \Delta$
\DisplayProof
\]
विगमन गृहीतकानुसार
$B, \Gamma \fCenter \Delta$
ची सिद्धता~$\pi_1''$ मिळते.
आता \CutCS{} ने संपणारी
पुढील सिद्धता पाहा:
\[
\AxiomC{}
\RightLabel{$\pi_2'$}
\Deduce$\Gamma \fCenter \Delta, B$
\AxiomC{}
\RightLabel{$\pi_1''$}
\Deduce$B, \Gamma \fCenter \Delta$
\RightLabel{\CutCS}
\BinaryInf$\Gamma \fCenter \Delta$
\DisplayProof
\]
तिची कट-कोटीही $\pi$
च्या कट-कोटीपेक्षा कमी आहे.
म्हणून विगमन गृहीतकानुसार
$\Gamma \fCenter \Delta$
ची \Log{G3c}-सिद्धता~$\pi'$
मिळते.
\item $A \ident \lforall[x][B(x)]$
असेल, तर $\pi$ चा शेवट असा:
\[
\AxiomC{}
\RightLabel{$\pi_1'$}
\Deduce$\Gamma \fCenter \Delta, B(c)$
\RightLabel{\RightR{\lforall}}
\UnaryInf$\Gamma \fCenter \Delta, \lforall[x][B(x)]$
\LeftSubproofLabel{$\pi_1$}
\AxiomC{}
\RightLabel{$\pi_2'$}
\Deduce$B(t), \lforall[x][B(x)], \Gamma \fCenter \Delta$
\RightLabel{\LeftR{\lforall}}
\UnaryInf$\lforall[x][B(x)], \Gamma \fCenter \Delta$
\RightSubproofLabel{$\pi_2$}
\RightLabel{\CutCS}
\BinaryInf$\Gamma \fCenter \Delta$
\DisplayProof
\]
विगमन गृहीतकानुसार
पुढील सिद्धतेतून \CutCS{}
काढता येतो:
\[
\AxiomC{}
\RightLabel{$\pi_1''$}
\Deduce$B(t), \Gamma \fCenter \Delta, \lforall[x][B(x)]$
\AxiomC{}
\RightLabel{$\pi_2'$}
\Deduce$B(t), \lforall[x][B(x)], \Gamma \fCenter \Delta$
\RightLabel{\CutCS}
\BinaryInf$B(t), \Gamma \fCenter \Delta$
\DisplayProof
\]
येथे $\pi_1''$ ही $\pi_1$
पासून (फक्त~$\pi_1'$ पासून
नव्हे!) उंची न वाढवणारे
दुर्बलीकरण लावून मिळते.
(या सिद्धता ची कट-कोटी
तेवढीच, पण कटची उंची
कमी आहे.) त्यामुळे
$B(t), \Gamma \fCenter \Delta$
ची सिद्धता~$\pi_2''$ मिळते.

सिद्धता~$\pi_1'$ ची
अंतिम क्रमवर्ती
$\Gamma \Sequent \Delta, B(c)$
आहे, जेथे $c$ हा
$\RightR{\lforall}$ अनुमानाचा
आयगेन चर आहे. म्हणून $c$ हे
$B(x)$, $\Gamma$ किंवा~$\Delta$
यांत येत नाही. सिद्धता~$\pi$
नियमित आहे असे मानू; त्यामुळे
$c$ हा फक्त पुढील
$\RightR{\lforall}$ अनुमानाचा
आयगेन चर आहे आणि फक्त त्याच्या
वर येतो. म्हणजे तो फक्त~$\pi_1'$
मध्ये येतो आणि $\pi_1'$ मधील
कोणत्याही अनुमानाचा आयगेन चर
नसतो. $c$ हे $\pi_2$ मध्ये
येत नसल्यामुळे $t$ मध्येही
येत नाही.
\ref{pt:seq:qua:lem:subst-regular}
नुसार सिद्धता~$\Subst{\pi_1'}{t}{c}$
ही $\Gamma \Sequent \Delta, B(t)$
ची सिद्धता आहे. आता पुढील
सिद्धता ला विगमन गृहीतक लावू:
\[
\AxiomC{}
\RightLabel{$\Subst{\pi_1'}{t}{c}$}
\Deduce$\Gamma \fCenter \Delta, B(t)$
\AxiomC{}
\RightLabel{$\pi_2''$}
\Deduce$B(t), \Gamma \fCenter \Delta$
\RightLabel{\CutCS}
\BinaryInf$\Gamma \fCenter \Delta$
\DisplayProof
\]
(कटचे सूत्र~$B(t)$
असल्यामुळे विगमन गृहीतक
लागू होते; या सिद्धता ची
कट-कोटी $\pi$ पेक्षा कमी आहे.)
\item $A \ident \lexists[x][B(x)]$
हा प्रसंग सरावासाठी ठेवतो.
\end{enumerate}
\end{proof}

\begin{prob}
\ref{pt:cut:top:lem:cut-adm-G3c} च्या
सिद्धतेतील प्रसंग~D चा
पुढील उपप्रसंग तपासा:
$\pi_1$ चा शेवट
$\LeftR{\lif}$ ने होतो,
पण त्याच्या अंतिम अनुमानात
$A$ मुख्य नसते.
(टीप: येथे नेमके काय सिद्ध
करायचे, म्हणजे या प्रसंगात
$\pi$ चे रूप काय, हे स्पष्ट
मांडणेही सरावाचा भाग आहे.)
\end{prob}

\begin{prob}
\ref{pt:cut:top:lem:cut-adm-G3c} च्या
सिद्धतेतील प्रसंग~D चा
पुढील उपप्रसंग तपासा:
$\pi_1$ चा शेवट
$\LeftR{\lforall}$ ने होतो,
पण त्याच्या अंतिम अनुमानात
$A$ मुख्य नसते.
\end{prob}

\begin{prob}
\ref{pt:cut:top:lem:cut-adm-G3c} च्या
सिद्धतेतील प्रसंग~E चा
पुढील उपप्रसंग तपासा:
$\pi_1$ आणि $\pi_2$
या दोन्हींच्या अंतिम अनुमानात
$A$ मुख्य आहे आणि
$A \ident (B \land C)$.
\end{prob}

\begin{prob}
\ref{pt:cut:top:lem:cut-adm-G3c} च्या
सिद्धतेतील प्रसंग~E चा
पुढील उपप्रसंग तपासा:
$\pi_1$ आणि $\pi_2$
या दोन्हींच्या अंतिम अनुमानात
$A$ मुख्य आहे आणि
$A \ident \lexists[x][B(x)]$.
\end{prob}

प्रसंग~E मध्ये विगमन
गृहीतक ज्या \CutCS{} ने
संपणाऱ्या सिद्धता वर
लावतो, तिची कटची उंची
मूळ सिद्धतेच्या कटच्या
उंची पेक्षा मोठी
असू शकते. उदाहरणार्थ,
$A \ident (B \lif C)$
असताना $\pi$ चे दोन
\CutCS{} अनुमाने असणाऱ्या
सिद्धता मध्ये रूपांतर
केले:
\[
\AxiomC{}
\RightLabel{$\pi_2'$}
\Deduce$\Gamma \fCenter \Delta, B$
\AxiomC{}
\RightLabel{$\pi_1'$}
\Deduce$B, \Gamma \fCenter \Delta, C$
\AxiomC{}
\RightLabel{$\pi_2'''$}
\Deduce$C, B, \Gamma \fCenter \Delta$
\RightLabel{\CutCS}
\BinaryInf$B, \Gamma \fCenter \Delta$
\RightSubproofLabel{$\pi_1''$}
\RightLabel{\CutCS}
\BinaryInf$\Gamma \fCenter \Delta$
\DisplayProof
\]
वरचा \CutCS{}, जो
$\pi_1'$ आणि $\pi_2''$
यांच्यामध्ये आहे, त्याची
कट-कोटी आणि कटची उंची
दोन्ही $\pi$ पेक्षा कमी आहेत.
पण विगमन गृहीतक लावल्याने
मिळालेली सिद्धता~$\pi_1''$
हिची कटची उंची आता
$\pi$ पेक्षा कमी असण्याची
गरज नाही. तरी खालच्या
\CutCS{} चे कटचे सूत्र
$B$ असल्यामुळे त्याची
\emph{कोटी} $\pi$ च्या
कट-कोटीपेक्षा कमी आहे.
येथे $\pi_2'''$ हे आधीप्रमाणे $\pi_2''$ चे
दुर्बलीकरण आहे; त्यामुळे वरच्या कटाच्या दोन्ही आधारांत
$B,\Gamma$ हा सामायिक संदर्भ आहे. खालचा कट $B$
काढतो; त्याचा निष्कर्ष $\Gamma\Sequent\Delta$ आहे.

\ref{pt:cut:top:lem:cut-adm-G3c} नुसार
सिद्धता मधील एक सर्वोच्च
\CutCS{} अनुमान काढता येते.
हे उपप्रमेय सर्वोच्च~\CutCS{}
ने संपणाऱ्या उप-सिद्धता
वर एकामागोमाग लावून
सिद्धता मधील कितीही
\CutCS{} अनुमाने काढता
येतात.

\begin{thm}\label{pt:cut:top:thm:cut-elim-G3c-topmost}
$\Log{G3c} + \CutCS$ मधील
कोणतीही सिद्धता $\pi$
ही~\Log{G3c} मधील
सिद्धता मध्ये रूपांतरित
करता येते.
\end{thm}

\begin{proof}
  $\pi$ मधील \CutCS{}
  अनुमानांची संख्या~$n$
  हिच्यावर विगमन करू.
  $n=0$ असेल, तर काही
  करायचे नाही: $\pi$
  ही आधीच~\Log{G3c}
  मधील सिद्धता आहे.
  अन्यथा, सर्वोच्च
  \CutCS{} अनुमानावर
  संपणारी उप-सिद्धता~$\pi'$
  विचारात घ्या. तिचे
  एकमेव \CutCS{} अंतिम
  अनुमान असते.
  \ref{pt:cut:top:lem:cut-adm-G3c}
  लावून तिचे \CutCS{}
  विरहित सिद्धता~$\pi''$
  मध्ये रूपांतर करा आणि
  मूळ सिद्धतेत
  उप-सिद्धता~$\pi'$
  च्या जागी~$\pi''$ ठेवा.
  मिळालेल्या सिद्धता
  मध्ये $n-1$ \CutCS{}
  अनुमाने उरतात; म्हणून
  विगमन गृहीतक लागू होते.
\end{proof}













\section{सर्वोच्च कोटीचे कट काढून टाकणे}\label{pt:cut:inv:sec}

\ref{pt:seq:inv:lem:G3c-invert} च्या सिद्धतेत दिसले
की \Log{G3c} ने एखाद्या तार्किक नियमाचा निष्कर्ष
(तो नियम \RightR{\lexists} किंवा \LeftR{\lforall}
नसेल तर) सिद्ध केला, तर त्या नियमाची प्रत्येक
आधार-क्रमवर्तीसुद्धा सिद्धतेची उंची न वाढवता
सिद्ध करता येते. याहून अधिक दाखवता येते:
हेच $\Log{G3c} + \CutCS$ साठीही लागू होते.

आधीप्रमाणे, \CutCS{} अनुमानाची \emph{कट-कोटी}
म्हणजे त्यातील कटच्या सूत्राची खोली.

\begin{lem}\label{pt:cut:inv:lem:inv-G3c-cut}
$\Log{G3c} + \CutCS$ मध्ये \RightR{\lexists}
किंवा \LeftR{\lforall} सोडून अन्य तार्किक नियम~$R$
चा निष्कर्ष सिद्धता~$\pi$ ने सिद्ध होत असेल,
तर त्या नियमाची प्रत्येक आधार-क्रमवर्ती
सिद्धता~$\pi'$ ने (किंवा नियमाला दोन आधार
असल्यास सिद्धता~$\pi'$, $\pi''$ यांनी)
सिद्ध करता येते. शिवाय,
(a)~$\pheight{\pi'}\le\pheight{\pi}$ (आणि
$\pheight{\pi''}\le\pheight{\pi}$);
तसेच (b)~प्रत्येक कोटीसाठी $\pi'$ मधील
(आणि स्वतंत्रपणे $\pi''$ मधील) त्या कोटीच्या
\CutCS{} अनुमानांची संख्या $\pi$ मधील संख्येपेक्षा
मोठी नसते.
\end{lem}

\begin{proof}
पूर्वप्रमेय~\ref{pt:seq:inv:lem:G3c-invert} आणि~\ref{pt:seq:inv:lem:invert-quant}
यांच्या सिद्धता तपासल्यास हा दावा मिळतो. ती
सिद्धता~$\pheight{\pi}$ वरील विगमनाने होती.
आधार प्रसंगात एक स्वयंसिद्धक दुसऱ्या
स्वयंसिद्धकाने बदलले होते; म्हणून (a) आणि (b)
लगेच मिळतात. $\pi$ चा शेवटचा नियम~$R$
असेल, तर त्याच्या तत्काळ उप-सिद्धता या
आपल्याला हव्या असलेल्या सिद्धता~$\pi_i$
आहेत; त्यांची उंची लहान असते आणि प्रत्येक कोटीच्या
कटांची संख्या संपूर्ण $\pi$ मधील संख्येपेक्षा मोठी नसते.
दोन आधारांच्या नियमात एक आधार निवडल्यास दुसऱ्या
शाखेतील कट राहतीलच असे नाही; (b) साठी समानता नको,
ही वरची सीमाच आवश्यक आहे.
इतर प्रसंगांत $\pi$ च्या तत्काळ उप-सिद्धता वर,
म्हणजे शेवटच्या अनुमान~$R$ च्या आधारांकडे
नेणाऱ्या सिद्धतांवर, विगमन गृहीतक लावून
परिणामांवर तोच नियम $R$ लावला होता.
नव्या सिद्धतेच्या तत्काळ उप-सिद्धता साठी
(a) व (b) खऱ्या असल्याने संपूर्ण
सिद्धतेसाठीही त्या खऱ्या आहेत. शेवटचे
अनुमान~\CutCS{} असेल तो प्रसंग तपासणे
उरते. पुन्हा फक्त एक उदाहरण पाहू:
$\pi$ ही \LeftR{\land} च्या
$B \land C, \Gamma \Sequent \Delta$
या निष्कर्षाची सिद्धता असून
\CutCS{} ने संपते असे समजा:
\[
\AxiomC{}
\RightLabel{$\pi_1$}
\Deduce$B \land C, \Gamma \fCenter \Delta, A$
\AxiomC{}
\RightLabel{$\pi_2$}
\Deduce$A, B \land C, \Gamma \fCenter \Delta$
\RightLabel{\CutCS}
\BinaryInf$B \land C, \Gamma \fCenter \Delta$
\DisplayProof
\]
विगमन गृहीतक $\pi_1$ आणि~$\pi_2$
यांना लागू होते; त्या दोन्ही
$\LeftR{\land}$ च्या निष्कर्षाच्या
उदाहरणांच्या सिद्धता आहेत.
त्यातून अनुक्रमे $B, C, \Gamma \fCenter
\Delta, A$ आणि $A, B, C, \Gamma \fCenter \Delta$
यांच्या सिद्धता $\pi_1'$ आणि~$\pi_2'$
मिळतात. \CutCS{} वापरून त्यांना जोडल्यास
पुढील~$\pi'$ मिळते:
\[
\AxiomC{}
\RightLabel{$\pi_1'$}
\Deduce$B, C, \Gamma \fCenter \Delta, A$
\AxiomC{}
\RightLabel{$\pi_2'$}
\Deduce$A, B, C, \Gamma \fCenter \Delta$
\RightLabel{\CutCS}
\BinaryInf$B, C, \Gamma \fCenter \Delta$
\DisplayProof
\]
येथे (a) आणि (b) स्पष्टपणे पूर्ण होतात.
\end{proof}

\CutCS{} ऐवजी \Cut{} वापरला असता ही
पद्धत चालली नसती: शेवटच्या सिद्धता
च्या निष्कर्षातील पूर्वांगात
$B, C, \Gamma$ च्या दोन प्रती आणि
उत्तरांगात $\Delta$ च्या दोन प्रती आल्या
असत्या. मग आकुंचनाची ग्राह्यता वापरावी
लागली असती; पण
\ref{pt:seq:inv:prop:G3c-cont-adm}
ची सिद्धता या उलटवण्याच्या उपप्रमेयावरच
आधारित आहे. त्यामुळे युक्तिवाद चक्रीय ठरला असता.
या उलटवण्याच्या उपप्रमेयानंतर दुर्बलीकरण आणि आकुंचनाची
ग्राह्यताही $\Log{G3c}+\CutCS$ पर्यंत वाढवता येते;
त्यामुळे पुढील सिद्धतेत कमी कोटीचे कट असतानाही आकुंचन
वापरता येते. आधीच्या उंचीवरील विगमनात एवढाच नवा प्रसंग
आहे: अंतिम नियम $\CutCS$ असेल, तर दोन्ही आधारांतील
सामायिक संदर्भावर विगमन गृहीतक लावून तोच कट पुन्हा
जोडायचा. कट-सूत्र आणि त्याची कोटी बदलत नाहीत. इतर
नियमांसाठी आधीची सिद्धता वापरायची; आकुंचनाच्या मुख्य
प्रसंगांत आता वरील विस्तारित उलटवणे वापरायचे. त्यातून
उंची आणि कटांच्या सर्वोच्च कोटीत वाढ होत नाही.
संख्यक नियमांच्या आयगेन चरांची नावे आवश्यकतेनुसार आधी
बदलून नव्या संदर्भाच्या चरांपासून वेगळी ठेवायची.

पूर्वप्रमेय~\ref{pt:cut:inv:lem:inv-G3c-cut} वापरून
कट-निर्मूलनाची पर्यायी सिद्धता देता येते.
या सिद्धतेत आपण सर्वोच्च स्थानी असलेली
\CutCS{} अनुमाने एकामागोमाग काढत नाही;
त्याऐवजी सर्वोच्च कोटीची \CutCS{}
अनुमाने काढतो. म्हणजे
$\Log{G3c} + \CutCS$ मधील सिद्धता~$\pi$
घेऊन तिच्यातील कोणताही कट काढतो
(कदाचित त्याजागी कमी कोटीचे कट येतील)
आणि एकही कट उरेपर्यंत असेच करत राहतो.
अट एवढीच की निवडलेल्या कटच्या वर
त्याच किंवा त्याहून मोठ्या कोटीचे
कट नसावेत.

\begin{lem}\label{pt:cut:inv:lem:max-cut-red-G3c}
$\Log{G3c}+\CutCS$ मधील सिद्धता~$\pi$
ही $n$ कोटीच्या \CutCS{} अनुमानाने
संपत असेल आणि इतरत्र फक्त
$\Log{G3c}$ चे नियम व $<n$ कोटीची
\CutCS{} अनुमाने वापरत असेल, तर
$\Log{G3c} + \CutCS$ मध्ये
त्याच अंतिम क्रमवर्तीची अशी सिद्धता~$\pi'$
असते की तिच्यातील प्रत्येक \CutCS{}
अनुमानाची कोटी~$<n$ असते.
\end{lem}

\begin{proof}
सिद्धता~$\pi$ चा शेवट पुढीलप्रमाणे असतो:
\[
\AxiomC{}
\RightLabel{$\pi_1$}
\Deduce$\Gamma \fCenter \Delta, A$
\AxiomC{}
\RightLabel{$\pi_2$}
\Deduce$A, \Gamma \fCenter \Delta$
\RightLabel{\CutCS}
\BinaryInf$\Gamma \fCenter \Delta$
\DisplayProof
\]
आता~$A$ च्या रूपानुसार प्रसंग वेगळे करू.

$A$ हे अणुसूत्र आहे असे समजा. $\pi_1$
आणि $\pi_2$ मधील सर्व \CutCS{} अनुमानांची कोटी
$< \depth{A} = 0$ असली पाहिजे; म्हणून त्यांपैकी
कोणत्याही सिद्धतेत \CutCS{} अनुमान असू शकत नाही.
म्हणून $\pi_1$ आणि $\pi_2$ दोन्ही कटविरहित आहेत.
$\pheight{\pi_2}$ वरील विगमनाने
$\Gamma \Sequent \Delta$ ची कटविरहित सिद्धता
आहे हे दाखवू. $\pheight{\pi_2} = 0$ असेल,
तर $A, \Gamma \Sequent \Delta$ हे स्वयंसिद्धक आहे.
$\lfalse$ हे $\Gamma$ मध्ये असेल, तर निष्कर्ष स्वतःच
असत्य-स्वयंसिद्धक आहे. अन्यथा $A$ मुख्य नसेल, तर
काही अणुसूत्र $C$ हे $\Gamma$ आणि $\Delta$ या दोन्हींत
घटक असते; म्हणून
$\Gamma \Sequent \Delta$ हे स्वतःच स्वयंसिद्धक
आहे. $A\ident\lfalse$ असल्याचा प्रसंग
\ref{pt:cut:top:lem:cut-adm-G3c} च्या सिद्धतेतील प्रसंग B प्रमाणे
हाताळता येतो: कटविरहित $\pi_1$ च्या उत्तरांगातील कटासाठीची
असत्याची प्रत काढल्यास हवी ती सिद्धता मिळते.
उरलेल्या मुख्य अणुसूत्राच्या प्रसंगात
$\Delta = \Delta', A$. या प्रसंगात $\pi_1$
चा शेवट $\Gamma \Sequent \Delta', A, A$
असा असतो. \ref{pt:seq:inv:prop:G3c-cont-adm},
म्हणजे आकुंचनाची ग्राह्यता, लावल्यास
$\Gamma \Sequent \Delta', A$ ची
कटविरहित सिद्धता मिळते.
$\pheight{\pi_2} > 0$ असेल, तर
सिद्धता नियम~$R$ ने संपते. त्या नियमाची
एक आधार-क्रमवर्ती घ्या: तिचे रूप
$A, \Gamma', \Pi \Sequent \Delta', \Lambda$
असते, जेथे $\Pi$ व $\Lambda$ ही
$R$ ची सक्रिय सूत्रे आहेत.
\ref{pt:seq:adm:prop:weak-G3c-adm}
हे $\pi_1$ आणि निवडलेल्या आधारावर संपणाऱ्या
$\pi_2$ च्या तत्काळ उपसिद्धतेला लावल्यास
$\Gamma, \Gamma', \Pi \Sequent
\Delta, \Delta', \Lambda, A$ आणि
$A, \Gamma, \Gamma', \Pi \Sequent
\Delta, \Delta', \Lambda$ यांच्या
सिद्धता मिळतात. दुसरी सिद्धता $\pi_2$ पेक्षा
कमी उंचीची राहते; म्हणून त्यांना विगमन गृहीतक लागू होते. त्यामुळे~$R$ च्या प्रत्येक
आधारासाठी $\Gamma, \Gamma', \Pi \Sequent
\Delta, \Delta', \Lambda$ ची सिद्धता
मिळते. $R$ लावल्यास
$\Gamma, \Gamma \Sequent \Delta, \Delta$
मिळते; आणि \ref{pt:seq:inv:prop:G3c-cont-adm}
मुळे $\Gamma \Sequent \Delta$ ची
कटविरहित सिद्धता मिळते.


$A$ अणुसूत्र नसेल, तर त्याच्या मुख्य संकारक नुसार प्रसंग पाहू. येथे $A$ हाच
कट-सूत्र आहे. प्रत्येक प्रसंगात
\ref{pt:cut:inv:lem:inv-G3c-cut} वापरून $A$ मुख्य सूत्र
असलेल्या डाव्या व उजव्या नियमांच्या आधारांच्या
सिद्धता मिळवू. त्यानंतर
\ref{pt:cut:top:lem:cut-adm-G3c} च्या सिद्धतेतील
प्रसंग~E प्रमाणे पुढे जाऊ.
\begin{enumerate}
  \item $A \ident \lnot B$. सिद्धता $\pi_1$
  आणि $\pi_2$ यांचे शेवट अनुक्रमे
  $\Gamma \Sequent \Delta, \lnot B$ आणि
  $\lnot B, \Gamma \Sequent \Delta$ असे आहेत.
  \ref{pt:cut:inv:lem:inv-G3c-cut} नुसार
  $B, \Gamma \Sequent \Delta$ आणि
  $\Gamma \Sequent \Delta, B$ यांच्या
  सिद्धता $\pi_1'$ आणि $\pi_2'$ मिळतात.
  सिद्धता~$\pi'$ पुढीलप्रमाणे आहे:
  \[
  \AxiomC{}
  \RightLabel{$\pi_2'$}
  \Deduce$\Gamma \fCenter \Delta, B$
  \AxiomC{}
  \RightLabel{$\pi_1'$}
  \Deduce$B, \Gamma \fCenter \Delta$
  \RightLabel{\CutCS}
  \BinaryInf$\Gamma \fCenter \Delta$
  \DisplayProof
  \]
  \item $A \ident (B \lor C)$.
  सिद्धता $\pi_1$ आणि $\pi_2$
  यांचे शेवट
  $\Gamma \Sequent \Delta, B \lor C$ आणि
  $B \lor C, \Gamma \Sequent \Delta$ असे आहेत.
  $\pi_1$ ला \RightR{\lor} चे उलटवणे
  लावून \ref{pt:cut:inv:lem:inv-G3c-cut} नुसार
  $\Gamma \Sequent \Delta, B, C$
  ची सिद्धता $\pi_1'$ मिळते.
  $\pi_2$ ला \LeftR{\lor} चे उलटवणे
  लावून \ref{pt:cut:inv:lem:inv-G3c-cut} नुसार
  $B, \Gamma \Sequent \Delta$ ची
  सिद्धता $\pi_2'$ आणि
  $C, \Gamma \Sequent \Delta$ ची
  सिद्धता $\pi_2''$ मिळते.
  \ref{pt:seq:adm:prop:weak-G3c-adm}
  हे $\pi_2''$ ला लावल्यास
  $C, \Gamma \Sequent \Delta, B$
  ची सिद्धता~$\pi_2'''$ मिळते.
  सिद्धता~$\pi'$ पुढीलप्रमाणे आहे:
  \[
  \AxiomC{}
  \RightLabel{$\pi_1'$}
  \Deduce$\Gamma \fCenter \Delta, B, C$
  \AxiomC{}
  \RightLabel{$\pi_2'''$}
  \Deduce$C, \Gamma \fCenter \Delta, B$
  \RightLabel{\CutCS}
  \BinaryInf$\Gamma \fCenter \Delta, B$
  \AxiomC{}
  \RightLabel{$\pi_2'$}
  \Deduce$B, \Gamma \fCenter \Delta$
  \RightLabel{\CutCS}
  \BinaryInf$\Gamma \fCenter \Delta$
  \DisplayProof
  \]
  \item $A \ident (B \land C)$.
  हा प्रसंग सरावासाठी आहे.
  \item $A \ident (B \lif C)$.
  सिद्धता $\pi_1$ आणि $\pi_2$
  यांचे शेवट
  $\Gamma \Sequent \Delta, B \lif C$ आणि
  $B \lif C, \Gamma \Sequent \Delta$ असे आहेत.
  $\pi_1$ ला \RightR{\lif} चे उलटवणे
  लावून \ref{pt:cut:inv:lem:inv-G3c-cut} नुसार
  $B, \Gamma \Sequent \Delta, C$
  ची सिद्धता $\pi_1'$ मिळते.
  $\pi_2$ ला \LeftR{\lif} चे उलटवणे
  लावून \ref{pt:cut:inv:lem:inv-G3c-cut} नुसार
  $\Gamma \Sequent \Delta, B$ ची
  सिद्धता $\pi_2'$ आणि
  $C, \Gamma \Sequent \Delta$ ची
  सिद्धता $\pi_2''$ मिळते.
  \ref{pt:seq:adm:prop:weak-G3c-adm}
  हे $\pi_2''$ ला लावल्यास
  $C, B, \Gamma \Sequent \Delta$
  ची सिद्धता~$\pi_2'''$ मिळते.
  सिद्धता~$\pi'$ पुढीलप्रमाणे आहे:
  \[
  \AxiomC{}
  \RightLabel{$\pi_2'$}
  \Deduce$\Gamma \fCenter \Delta, B$
  \AxiomC{}
  \RightLabel{$\pi_1'$}
  \Deduce$B, \Gamma \fCenter \Delta, C$
  \AxiomC{}
  \RightLabel{$\pi_2'''$}
  \Deduce$C, B, \Gamma \fCenter \Delta$
  \RightLabel{\CutCS}
  \BinaryInf$B, \Gamma \fCenter \Delta$
  \RightLabel{\CutCS}
  \BinaryInf$\Gamma \fCenter \Delta$
  \DisplayProof
  \]
  \item $A \ident \lforall[x][B(x)]$.
  सिद्धता $\pi_1$ आणि $\pi_2$
  यांचे शेवट
  $\Gamma \Sequent \Delta, \lforall[x][B(x)]$
  आणि $\lforall[x][B(x)], \Gamma \Sequent \Delta$
  असे आहेत. \ref{pt:cut:inv:lem:inv-G3c-cut} नुसार
  $\Gamma \Sequent \Delta, B(c)$ ची
  सिद्धता~$\pi_1'$ मिळते.

आता सिद्धता~$\pi_2$ पाहू. तिचा शेवट
$\lforall[x][B(x)], \Gamma \Sequent \Delta$
असा आहे. $\pi_2$ च्या अंतिम क्रमवर्तीपासून
वर जाताना, डाव्या बाजूला आलेली
$\lforall[x][B(x)]$ ची अशी प्रत्येक आस्थिती
चिन्हांकित करा की ती अंतिम क्रमवर्तीतील
$\lforall[x][B(x)]$ च्या आस्थितीपर्यंत
“पोहोचते”. म्हणजे:
(a)~$\pi_2$ च्या अंतिम क्रमवर्तीतील
$\lforall[x][B(x)]$ चिन्हांकित करा.
(b)~एखाद्या नियमाच्या निष्कर्षाच्या
संदर्भात $\lforall[x][B(x)]$ चिन्हांकित
असेल, तर त्या नियमाच्या आधारात किंवा
आधारांत तिच्याशी अनुरूप आस्थिती चिन्हांकित
करा. (c)~\LeftR{\lforall} नियमाच्या
निष्कर्षात $\lforall[x][B(x)]$ हे मुख्य
सूत्र असून चिन्हांकित असेल, तर त्याच्या
आधारातील $\lforall[x][B(x)]$ आणि
$B(t)$ यांच्या अनुरूप सक्रिय आस्थिती
चिन्हांकित करा. हा मागोवा $\pi_2$
च्या सिद्धतावृक्षावरच घेतला आहे.

आता $\lforall[x][B(x)]$ च्या
या सर्व चिन्हांकित आस्थिती पाहू. त्यांपैकी
कोणतीही \LeftR{\lforall} व्यतिरिक्त
अन्य नियमात सक्रिय नसते
(फक्त \LeftR{\lforall} ची सक्रिय
सूत्रे चिन्हांकित होतात).
$\pi_2$ मधील $\lforall[x][B(x)]$
च्या सर्व चिन्हांकित आस्थिती काढून टाका.
यातून योग्य सिद्धता मिळाली, तर तिचा शेवट
$\Gamma \Sequent \Delta$ असा आहे आणि
ती $\pi$ च्या जागी ठेवता येते.
याने सर्वोच्च कोटीचा एक कट काढला.

अन्यथा मिळालेली सिद्धता योग्य नाही:
काही \LeftR{\lforall} अनुमानांत सक्रिय
असलेली $\lforall[x][B(x)]$ ही
सूत्रे आपण काढून टाकली आहेत.
त्यामुळे पुढील रूपाची “अनुमाने” मिळाली:
\[
  \AxiomC{}
  \RightLabel{$\pi_t$}
  \Deduce$B(t), \Pi \fCenter \Lambda$
  \RightLabel{\LeftR{\lforall}}
  \UnaryInf$\Pi \fCenter \Lambda$
  \DisplayProof
  \]
अशा प्रत्येक सर्वोच्च “अनुमाना”च्या जागी
पुढील सिद्धता ठेवा:
\[
  \AxiomC{}
  \RightLabel{$\Subst{\pi_1'}{t}{c}[\Pi \Sequent \Lambda]$}
  \Deduce$\Gamma, \Pi \fCenter \Delta, \Lambda, B(t)$
  \AxiomC{}
  \RightLabel{$\pi_t[\Gamma \Sequent \Delta]$}
  \Deduce$B(t), \Gamma, \Pi \fCenter \Delta, \Lambda$
  \RightLabel{\CutCS}
  \BinaryInf$\Gamma, \Pi \fCenter \Delta, \Lambda$
  \DisplayProof
  \]
तिच्या खालील प्रत्येक क्रमवर्तीच्या
डावीकडे $\Gamma$ आणि उजवीकडे
$\Delta$ घाला. आधारात चिन्हांकित
$B(t)$ असलेली सर्व
\LeftR{\lforall} “अनुमाने”
काही~$B(t)$ वरील \CutCS{}
अनुमानाने बदलेपर्यंत हे पुन्हा करा.
मिळालेल्या सिद्धता ची अंतिम क्रमवर्ती
(आता ती खरोखर योग्य सिद्धता आहे)
$\Gamma, \dots, \Gamma \Sequent
\Delta, \dots, \Delta$ या रूपाची असते.
आकुंचनाची ग्राह्यता लावून तिचे
$\Gamma \Sequent \Delta$ वरील
सिद्धता मध्ये रूपांतर करा आणि
ती $\pi$ च्या जागी ठेवा. आपण
$\lforall[x][B(x)]$ वरील एक कट
$B(t)$ वरील (कदाचित अनेक) कटांनी
बदलला आहे; पण या सर्व कटांची कोटी
कमी आहे. $\pi_1$ किंवा $\pi_2$ मधील उरलेल्या किंवा प्रत करून
वापरलेल्या कटांची कोटीही कमीच आहे. पुढील आकुंचनासाठी
वर सिद्ध केलेली, सर्वोच्च कट-कोटी न वाढवणारी विस्तारित
ग्राह्यता वापरली आहे.
\end{enumerate}
\end{proof}

\begin{prob}
\ref{pt:cut:inv:lem:max-cut-red-G3c} च्या सिद्धतेतील
$A \ident (B \land C)$ हा प्रसंग तपासा.
म्हणजे, दोन उपसिद्धतांसह पुढीलप्रमाणे संपणारी
सिद्धता~$\pi$ असेल:
\[
\AxiomC{}
\RightLabel{$\pi_1$}
\Deduce$\Gamma \fCenter \Delta, B \land C$
\AxiomC{}
\RightLabel{$\pi_2$}
\Deduce$B \land C, \Gamma \fCenter \Delta$
\RightLabel{\CutCS}
\BinaryInf$\Gamma \fCenter \Delta$
\DisplayProof
\]
आणि $\pi_1$ व $\pi_2$ यांतील
सर्व कटांची कोटी $<\depth{B
\land C}$ असेल, तर
$\Gamma \Sequent \Delta$ ची
फक्त $<\depth{B \land C}$
कोटीचे कट असलेली सिद्धता~$\pi'$
असते हे दाखवा.
\end{prob}






\ref{pt:cut:inv:lem:max-cut-red-G3c} नुसार
सिद्धता मधील सर्वोच्च कोटीचे
\CutCS{} अनुमान कमी कोटीच्या
\CutCS{} अनुमानांनी बदलता येते.
हे उपप्रमेय पुन्हा वापरून,
सर्वोच्च~\CutCS{} ने संपणाऱ्या
सर्वोच्च उप-सिद्धता वर ते एकामागोमाग
लावल्यास कोणत्याही सिद्धता मधून कितीही
\CutCS{} अनुमाने काढता येतात.


\begin{thm}\label{pt:cut:inv:thm:cut-elim-G3c-largest}
$\Log{G3c} + \CutCS$ मधील कोणतीही
सिद्धता $\pi$ ही~\Log{G3c} मधील
सिद्धतेत रूपांतरित करता येते.
\end{thm}

\begin{proof}
  प्रथम $\pi$ मधील सर्वोच्च कोटीचे
  सर्व कट काढता येतात हे सिद्ध करू.
  $d$ ही $\pi$ मधील कटांची सर्वोच्च
  कोटी आणि $n$ ही $d$ कोटीच्या कटांची
  संख्या असू द्या. $n$ वरील विगमनाने
  सिद्धता करू. $n=0$ असेल, तर
  सिद्ध करायचे काही नाही. $n>0$
  असेल, तर $d$ कोटीचा सर्वोच्च कट
  निवडा आणि त्यावर संपणारी उप-सिद्धता
  $\pi_1$ असू द्या.
  \ref{pt:cut:inv:lem:max-cut-red-G3c} नुसार
  त्याच अंतिम क्रमवर्तीची, फक्त
  $<d$ कोटीचे कट असलेली,
  सिद्धता~$\pi_1'$ आहे.
  $\pi$ मध्ये $\pi_1$ च्या जागी
  $\pi_1'$ ठेवा. मिळालेल्या
  सिद्धता मध्ये $d$ कोटीचे
  $n-1$ कट असतात; म्हणून
  विगमन गृहीतक लागू होते.

आता $d$ वरील विगमनाने निष्कर्ष
मिळतो. $d=0$ असेल, तर सर्व कट
अणुसूत्रांवरील आहेत आणि मागील
निष्कर्षाने ते सर्व काढता येतात.
$d>0$ असेल, तर मागील निष्कर्षाने
$d$ कोटीचे सर्व कट $<d$ कोटीच्या
कटांनी बदललेली सिद्धता मिळते.
$\pi$ मधील कटांची सर्वोच्च कोटी
$d$ असल्यामुळे नव्या सिद्धतेतील
कटांची सर्वोच्च कोटी $<d$ होते;
म्हणून विगमन गृहीतक लागू होते.
\end{proof}













\section{मध्य-क्रमवर्ती प्रमेय आणि एरब्राँचे प्रमेय}\label{pt:cut:mid:sec}

कट-निर्मूलन प्रमेयाचा एक महत्त्वाचा
परिणाम म्हणजे मध्य-क्रमवर्ती प्रमेय.
त्यात असे म्हटले आहे: सर्व सूत्र
अग्र-संख्यापक रूपातील असलेल्या
$\Gamma \Sequent \Delta$ ची
सिद्धता~$\pi$ असेल, तर अशी
सिद्धता~$\pi'$ असते की तिच्यात
$S \ident \Pi \Sequent
\Lambda$ ही क्रमवर्ती येते
(तिला \emph{मध्य-क्रमवर्ती} म्हणतात)
आणि~$\pi'$ मध्ये $S$ च्या
वरची सर्व अनुमाने विधानीय,
तर $S$ च्या खालची सर्व अनुमाने
संख्यापकांची असतात.
(सूत्र~$A$ हे
\emph{अग्र-संख्यापक} (किंवा
\emph{अग्र-संख्यापक रूपातील}) असते,
जर ते पुढील रूपात असेल:
\[\mathsf{Q}_1x_1\dots\mathsf{Q}_nx_n\,B(x_1,\dots,x_n)\]
येथे प्रत्येक~$\mathsf{Q}_i$
हा $\lforall$ किंवा~$\lexists$
असतो आणि $B$ संख्यापकविरहित असते.) मध्य-क्रमवर्ती प्रमेयाचा
एक महत्त्वाचा परिणाम म्हणजे
एरब्राँचे प्रमेय. त्याचे व्यापक
रूप थोडे कठीण आहे; साधारणतः
ते असे सांगते: प्रथम-क्रम
तर्कशास्त्रातील प्रत्येक
सिद्धतायोग्य वाक्य~$A$
साठी असे विधानीय
सर्वतः सत्य सूत्र~$A'$
असते की त्यावरून $A$
सिद्ध करता येते.
$B$ संख्यापकविरहित
असताना $\lexists[x][B(x)]$
या रूपातील वाक्ये साठी
पुढील विधान आहे:

\begin{thm}[एरब्राँचे प्रमेय]
$B(x)$ संख्यापकविरहित असू द्या
आणि $\Proves
\lexists[x][B(x)]$ असे समजा.
मग अशी पदे~$t_1$, \dots, $t_n$
असतात की
$\Proves B(t_1) \lor \dots \lor B(t_n)$.
\end{thm}

$B(t_1) \lor \dots \lor B(t_n)$
या सूत्राला
$\lexists[x][B(x)]$ चा
\emph{एरब्राँचा विकल्पयोग} म्हणतात.
$B$ संख्यापकविरहित असल्याने
हा विकल्पयोगही संख्यापकविरहित आहे.
म्हणून तो सिद्धतायोग्य असेल,
तर कटविरहित सिद्धतेने तो
सिद्धतायोग्य आहे; परिणामी,
केवळ विधानीय अनुमाने
वापरणारी सिद्धता आहे.
म्हणून तो सर्वतः सत्य आहे.

एरब्राँच्या प्रमेयातील कठीण
भाग म्हणजे प्रत्येक
सिद्धतायोग्य
सूत्र~$\lexists[x][B(x)]$
साठी एरब्राँचा विकल्पयोग
अस्तित्वात आहे हे दाखवणे.
उलट दिशा सोपी आहे:
$\lexists[x][B(x)]$ चा
एरब्राँचा विकल्पयोग असेल,
तर हे सूत्र सिद्धतायोग्य आहे.
कारण एरब्राँचा विकल्पयोग
$\lexists[x][B(x)]$
ला निष्पन्न करतो. उदाहरणार्थ,
$n=2$ असताना~$\Log{G3c}$
मधील ही सिद्धता पाहा:
\[
\Axiom$B(t_1) \fCenter \lexists[x][B(x)], B(t_1), B(t_2)$
\Axiom$B(t_2) \fCenter \lexists[x][B(x)], B(t_1), B(t_2)$
\RightLabel{\LeftR{\lor}}
\BinaryInf$B(t_1) \lor B(t_2) \fCenter \lexists[x][B(x)], B(t_1), B(t_2)$
\RightLabel{\RightR{\lexists}}
\UnaryInf$B(t_1) \lor B(t_2) \fCenter \lexists[x][B(x)], B(t_2)$
\RightLabel{\RightR{\lexists}}
\UnaryInf$B(t_1) \lor B(t_2) \fCenter \lexists[x][B(x)]$
\DisplayProof
\]

$\Sequent
\lexists[x][B(x)]$ ची
सिद्धता~$\pi$ दिली असता
त्यातून एरब्राँचा विकल्पयोग
मिळवण्यासाठी, कटविरहित
सिद्धता~$\pi$ मधील सर्व
$\RightR{\lexists}$
अनुमाने शेवटी सलग येतात
हा प्रसंग पाहा:
\[
\AxiomC{}
\RightLabel{$\pi_1$}
\Deduce$ \fCenter \lexists[x][B(x)], B(t_1), \dots, B(t_n)$
\RightLabel{\RightR{\lexists}}
\UnaryInf$ \fCenter \lexists[x][B(x)], B(t_2), \dots, B(t_n)$
\RightLabel{$(n-1) \times \RightR{\lexists}$}
\Deduce$ \fCenter \lexists[x][B(x)]$
\DisplayProof
\]
~$\pi$ मध्ये
$\lexists[x][B(x)]$
सक्रिय असलेली हीच सर्व
$\RightR{\lexists}$ अनुमाने
असतील, तर~$\pi_1$
मध्ये इतर संख्यापक अनुमाने
नसतात. कारण $B(x)$
मध्ये अन्य संख्यापक नाहीत
आणि अंतिम क्रमवर्तीच्या
डावीकडे सूत्रे नाहीत.
म्हणजे $ \fCenter \lexists[x][B(x)], B(t_1),
\dots, B(t_n)$ ही~$\pi$
ची मध्य-क्रमवर्ती आहे.
शिवाय मध्य-क्रमवर्तीच्या
वर संख्यापक अनुमाने
नसल्यामुळे तिच्या वरच्या
सर्व क्रमवर्तींत
$\lexists[x][B(x)]$
असते, पण ते सक्रिय किंवा
मुख्य नसते. म्हणून
~$\pi_1$ या सिद्धता
मधून ते काढून टाकता येते
आणि $\fCenter B(t_1), \dots, B(t_n)$
याची सिद्धता मिळते.
मग $\RightR{\lor}$ चा
$n-1$ वेळा उपयोग करून
एरब्राँच्या विकल्पयोगाची
सिद्धता,
$\Sequent B(t_1) \lor \dots \lor B(t_n)$,
मिळते.

अडचण अशी की
$\Sequent \lexists[x][B(x)]$
ची कटविरहित सिद्धता
असली तरी तिची सर्व
\RightR{\lexists} अनुमाने
शेवटी एका सलग गटात
येतात असे गृहीत धरता येत
नाही. त्या \RightR{\lexists}
अनुमानांमध्ये एखादे
$B(t_i)$ मुख्य असलेली
अनुमाने येऊ शकतात.
म्हणून मध्य-क्रमवर्ती
प्रमेयाची गरज आहे.

\begin{thm}[मध्य-क्रमवर्ती प्रमेय]\label{pt:cut:mid:thm:midsequent}
$\Gamma \Sequent \Delta$ ची
~\Log{G3c} मधील
सिद्धता~$\pi$ असेल आणि
$\Gamma$ व $\Delta$ मधील
सर्व सूत्रे
अग्र-संख्यापक असतील,
तर अशी सिद्धता~$\pi'$
असते की तिच्यात
$\Pi \Sequent \Lambda$
ही क्रमवर्ती येते;
तिच्या खालची सर्व
अनुमाने संख्यापकांची,
आणि वरची सर्व अनुमाने
विधानीय असतात.
\end{thm}

\begin{proof}
  ~$\pi$ मधील प्रत्येक संख्यापक
  अनुमानाचे \emph{क्रममान} म्हणजे
  त्याच्या खालच्या विधानीय
  अनुमानांची संख्या असे
  परिभाषित करू; आणि
  $\pi$ चे क्रममान~$o(\pi)$
  म्हणजे त्यातील सर्व
  संख्यापक अनुमानांच्या
  क्रममानांची बेरीज.
  $o(\pi) = 0$ असेल, तर
  कोणत्याही संख्यापक
  अनुमानाच्या खाली विधानीय
  अनुमान नसते आणि काम संपते:
  संख्यापक अनुमानच नसल्यास
  अंतिम क्रमवर्ती मध्य-क्रमवर्ती
  घेता येते. अन्यथा सर्व
  संख्यापक अनुमानांना एकच
  आधारक्रमवर्ती असल्याने
  एक सर्वोच्च संख्यापक
  अनुमान असते; त्याची
  आधारक्रमवर्ती मध्य-क्रमवर्ती असते.

  $o(\pi) > 0$ असेल, तर
  क्रममान $>0$ असलेले सर्वात
  खालचे संख्यापक अनुमान
  निवडा. त्याचे क्रममान $>0$
  असल्याने त्याच्या खाली
  विधानीय अनुमान असलेच पाहिजे.
  ते अशा अनुमानांपैकी
  सर्वात खालचे असल्याने
  त्याचा निष्कर्ष दुसऱ्या
  संख्यापक अनुमानाची
  आधारक्रमवर्ती असू शकत
  नाही: तसे असेल, तर
  त्या दुसऱ्या, अधिक
  खालच्या संख्यापक अनुमानाच्या
  खालीही विधानीय अनुमान
  आले असते. कोणते
  संख्यापक अनुमान आणि
  त्याखाली कोणते विधानीय
  अनुमान आहे, यानुसार
  (अनेक!) प्रसंग ठरतात.
  अंतिम क्रमवर्तीमध्ये
  केवळ अग्र-संख्यापक
  सूत्रे असल्यामुळे,
  संख्यापक अनुमानाचे
  मुख्य सूत्र पुढील
  विधानीय अनुमानात
  सक्रिय असू शकत नाही.
  अन्यथा त्या अनुमानाच्या
  मुख्य सूत्रात विधानीय
  संयोजकाच्या व्याप्तीत
  संख्यापक आला असता.
  उप-सूत्र गुणधर्मामुळे
  ते अंतिम क्रमवर्तीतील
  एखाद्या सूत्राचे
  उप-सूत्र असावे लागले
  असते. म्हणून संख्यापक
  अनुमानाचे मुख्य
  सूत्र हे खालच्या
  विधानीय अनुमानाच्या
  संदर्भातील घटक आहे.

  पुढील दोन उदाहरणे पाहू:

प्रथम $\RightR{\lforall}$
हे $\LeftR{\lif}$ च्या
उजव्या आधारक्रमवर्तीमध्ये
संपते असे समजा. मग
~$\pi$ ची संबंधित
उप-सिद्धता
ही उप-सिद्धता~$\pi_1$
मध्ये संपते:
\[
\AxiomC{}
\RightLabel{$\pi_2$}
\Deduce$\Gamma' \fCenter \Delta', \lforall[x][A(x)], B$
\AxiomC{}
\RightLabel{$\pi_3$}
\Deduce$C, \Gamma' \fCenter \Delta', A(c)$
\RightLabel{\RightR{\lforall}}
\UnaryInf$C, \Gamma' \fCenter \Delta', \lforall[x][A(x)]$
\RightLabel{\LeftR{\lif}}
\BinaryInf$B \lif C, \Gamma' \fCenter \Delta', \lforall[x][A(x)]$
\DisplayProof
\]
~$\pi$ नियमित आहे असे
मानता येते; म्हणून
आयगेन चर~$c$ हे~$\pi_3$
च्या बाहेर येत नाही.
विशेषतः ते $B
\lif C$ मध्ये किंवा
~$\pi_2$ मध्ये येत नाही.
आता सिद्धता~$\pi_1'$ पाहा:
\[
\AxiomC{}
\RightLabel{$\pi_2'$}
\Deduce$\Gamma' \fCenter \Delta', \lforall[x][A(x)], A(c), B$
\AxiomC{}
\RightLabel{$\pi_3$}
\Deduce$C, \Gamma' \fCenter \Delta', \lforall[x][A(x)], A(c)$
\RightLabel{\LeftR{\lif}}
\BinaryInf$B \lif C, \Gamma' \fCenter \Delta', \lforall[x][A(x)], A(c)$
\RightLabel{\RightR{\lforall}}
\UnaryInf$B \lif C, \Gamma' \fCenter \Delta', \lforall[x][A(x)], \lforall[x][A(x)]$
\DisplayProof
\]
येथे~$\pi_2'$ ही
~$\pi_2$ च्या उजवीकडे
~$A(c)$ जोडून दुर्बलीकरणाने
मिळते. (त्यामुळे नवीन
अनुमाने, विशेषतः क्रममानावर
परिणाम करू शकणारी
संख्यापक अनुमाने, जोडली जात
नाहीत. ~$\pi_2$ मध्ये
आयगेन चरे म्हणून
वापरलेले स्थिरांक~$A(c)$
मध्ये नसल्याने~$\pi_2$
मधील आयगेन चरांच्या
अटीही मोडत नाहीत.)
दुसऱ्या शाखेला आवश्यक
संदर्भ जोडण्यासाठीही
दुर्बलीकरण वापरता येते.
\RightR{\lforall} ची
आयगेन चराची अट अजूनही
पूर्ण आहे, कारण $c$
हे~$B \lif C$
मध्ये येत नाही.

आकुंचनाची ग्राह्यता वापरून
$B \lif C, \Gamma' \fCenter \Delta',
\lforall[x][A(x)]$ याची
सिद्धता~$\pi_1''$
मिळवतो. $\lforall[x][A(x)]$
साठी $\RightR{\Contraction}$
ची ग्राह्यता सिद्ध करताना
आणि त्या सिद्धतेत वापरलेल्या
$\RightR{\lforall}$ च्या
व्युत्क्रमणीयतेच्या सिद्धतेत
अनुमानांची क्रममाने वाढत नाहीत:
अनुमाने काढली जाऊ शकतात, म्हणून संबंधित
उरलेल्या अनुमानांच्या खालील विधानीय
अनुमानांची संख्या कमीही होऊ शकते.
म्हणून~$\pi_1''$ मध्ये
~$\pi_1$ पेक्षा अधिक
संख्यापक अनुमाने नाहीत,
आणि~$\pi_1''$ मधील
प्रत्येक संख्यापक अनुमानाच्या
खाली त्याच्या~$\pi_1$
मधील संबंधित अनुमानापेक्षा अधिक
विधानीय अनुमाने नाहीत.
अपवाद फक्त अंतिम
\RightR{\lforall} चा:
त्याच्या खाली~$\pi$
मध्ये \LeftR{\lif} होते,
ते आपण~$\pi_1''$ मध्ये
त्याच्या वर हलवले आहे.
~$\pi$ मधील $\pi_1$
च्या जागी $\pi_1''$ ठेवा.
मिळालेल्या सिद्धता
चे क्रममान कमी आहे, कारण
किमान $\RightR{\lforall}$
अनुमानाचे क्रममान (किमान)~$1$
ने कमी झाले. आता
विगमन गृहीतक लावा.

संख्यापक अनुमान
\RightR{\lexists} असल्याचे
प्रसंग आणखी सोपे आहेत:
आकुंचन लागत नाही आणि
आयगेन चराच्या अटीची
काळजीही करावी लागत नाही.
उदाहरणार्थ,
$\RightR{\lexists}$
नंतर $\RightR{\land}$
येतो असे समजा
(यावेळी डाव्या आधारक्रमवर्तीमध्ये).
संबंधित उप-सिद्धता
ही $\pi_1$ आहे:
\[
\AxiomC{}
\RightLabel{$\pi_2$}
\Deduce$\Gamma' \fCenter \Delta', \lexists[x][A(x)], A(t), B$
\RightLabel{\RightR{\lexists}}
\UnaryInf$\Gamma' \fCenter \Delta', \lexists[x][A(x)], B$
\AxiomC{}
\RightLabel{$\pi_3$}
\Deduce$\Gamma' \fCenter \Delta', \lexists[x][A(x)], C$
\RightLabel{\RightR{\land}}
\BinaryInf$\Gamma' \fCenter \Delta', \lexists[x][A(x)], B \land C$
\DisplayProof
\]
~$\pi$ मधील $\pi_1$
च्या जागी $\pi_1'$ ठेवा:
\[
\AxiomC{}
\RightLabel{$\pi_2$}
\Deduce$\Gamma' \fCenter \Delta', \lexists[x][A(x)], A(t), B$
\AxiomC{}
\RightLabel{$\pi_3'$}
\Deduce$\Gamma' \fCenter \Delta', \lexists[x][A(x)], A(t), C$
\RightLabel{\RightR{\land}}
\BinaryInf$\Gamma' \fCenter \Delta', \lexists[x][A(x)], A(t), B \land C$
\RightLabel{\RightR{\lexists}}
\UnaryInf$\Gamma' \fCenter \Delta', \lexists[x][A(x)], B \land C$
\DisplayProof
\]
येथे $\pi_3'$ ही
$\pi_3$ च्या उजवीकडे
~$A(t)$ जोडून
दुर्बलीकरणाने मिळते.
या दुर्बलीकरणात
~$\pi_3$ मधील
प्रत्येक क्रमवर्तीच्या
उजवीकडे केवळ $A(t)$
जोडले जाते; म्हणून
कोणत्याही संख्यापक
अनुमानाचे क्रममान बदलत
नाही. ~$\pi'$ आणि
~$\pi$ मधील
संख्यापक अनुमानांच्या
क्रममाने सारखीच आहेत;
फक्त~\RightR{\land}
च्या वर हलवलेल्या
\RightR{\lexists}
अनुमानाचे क्रममान~$1$
ने कमी होते. म्हणून
विगमन गृहीतक लागू होते.
\end{proof}

\begin{prob}
\ref{pt:cut:mid:thm:midsequent} च्या
सिद्धतेतील पुढील प्रसंग
वर्णन करा: \LeftR{\lexists}
हे~\RightR{\lif} च्या
आधारक्रमवर्तीमध्ये संपते;
आणि \LeftR{\lforall}
हे~\LeftR{\lor} च्या
डाव्या आधारक्रमवर्तीमध्ये
संपते.
\end{prob}

\begin{proof}[एरब्राँच्या प्रमेयाची सिद्धता]
$\pi$ चा शेवट
$\Sequent \lexists[x][A(x)]$
मध्ये होतो असे समजा.
\ref{pt:cut:mid:thm:midsequent} नुसार
$\pi$ चे मध्य-क्रमवर्ती~$S$
असलेल्या सिद्धता~$\pi'$
मध्ये रूपांतर करता येते;
तिचे रूप पुढीलप्रमाणे असले
पाहिजे: \[\Sequent
\lexists[x][A(x)], A(t_1), \dots, A(t_n).\]
$S$ च्या वरची सर्व
अनुमाने विधानीय असल्याने,
$\lexists[x][A(x)]$
हे $S$ च्या वरच्या
अनुमानात मुख्य असू शकत नाही.
ते अणुरूप नसल्याने
स्वयंसिद्धकातही मुख्य
असू शकत नाही. शिवाय,
$A(x)$ संख्यापकविरहित
आहे हे लक्षात घेतल्यास
$\lexists[x][A(x)]$
हे $A(t_1)$ चे
योग्य उपसूत्र असू शकत नाही;
म्हणून $S$ च्या वर
ते सक्रियही नसते.
त्यामुळे $S$ वर संपणाऱ्या
उप-सिद्धता मधून
$\lexists[x][A(x)]$
काढल्यावर
$\Sequent A(t_1), \dots, A(t_n)$
याची योग्य सिद्धता
मिळते. त्यातून
$\RightR{\lor}$ ची
$n-1$ अनुमाने लावून
$\Sequent A(t_1) \lor \dots \lor A(t_n)$
याची सिद्धता मिळते.
\end{proof}






















\section{माएहाराचे उपप्रमेय आणि क्रेगचे अंतर्वेशन प्रमेय}\label{pt:cut:itp:sec}

विल्यम क्रेग यांचे अंतर्वेशन प्रमेय असे सांगते:
$A \Entails B$ असेल, तर असे सूत्र~$C$
असते (त्याला \emph{अंतर्वेशक} म्हणतात) की
$A \Entails C$ आणि $C \Entails B$.
विशेष म्हणजे, $C$ मधील सर्व स्थिरांक आणि
विधेये हे $A$ व $B$ या दोन्हींमध्ये येतील
असे $C$ निवडता येते. (येथे कोणतीही फलने
नाहीत असे गृहीत धरतो.) ही अट नसेल, तर
$A$ आणि $B$ हेच सहजपणे अंतर्वेशक ठरतील.

अंतर्वेशन प्रमेय अभिजात प्रथम-क्रम तर्कशास्त्राला लागू होते.
त्याचा शोध प्रथम-क्रम तर्कशास्त्राचा “शेवटचा महत्त्वाचा गुणधर्म”
शोधला गेल्याप्रमाणे वर्णिला गेला आहे.
प्रतिमान उपपत्ती आणि स्वयंचलित तर्कपद्धती यांत त्याचे
महत्त्वाचे उपयोग आहेत. त्यामागील मूळ प्रेरणा आणि
उपयोग विज्ञानाच्या तत्त्वज्ञानात होता: एखाद्या उपपत्तीचे
स्वयंसिद्धकीकरण कोणत्या मूलभूत संज्ञांनी करता येते
किंवा येत नाही, हे तपासण्यासाठी ते वापरता येते.
यावरून बेथचे परिभाष्यता प्रमेयही निष्पन्न होते:
एखादी उपपत्ती संकल्पना अंतर्निहितपणे परिभाषित करू
शकत असेल, तर ती स्पष्टपणेही परिभाषित करू शकते.

गणिती तर्कशास्त्रातील अनेक निष्कर्षांप्रमाणे,
अंतर्वेशन प्रमेयाची सिद्धता प्रतिमान उपपत्तीच्या
तसेच सिद्धता उपपत्तीच्या पद्धतींनी करता येते.
सिद्धता-उपपत्तीय पद्धतीचा लाभ असा की ती
केवळ अंतर्वेशकाचे अस्तित्व दाखवत नाही, तर
$A \Entails B$ याच्या सिद्धता पासून
अंतर्वेशक मोजून देणारी अल्गोरिदमही देते;
उदाहरणार्थ, $A \Sequent B$ याची~\Log{G3c}
मधील सिद्धता वापरता येते.

$A$ मधील स्थिरांक आणि विधेये
यांचा संच $\Lang L(A)$ असू द्या आणि
$\Lang L(\Gamma) = \bigcup \Setabs{\Lang L(A)}{A \in
\Gamma}$ असू द्या. या संपूर्ण विभागात
फलने नसलेली सूत्रे आपण विचारात घेतो.

\begin{thm}[क्रेगचे अंतर्वेशन प्रमेय]\label{pt:cut:itp:thm:interpolation}
$A$ ची भाषा $\Lang L(A)$ आणि
$B$ ची $\Lang L(B)$ असू द्या.
$\Log{G3c} \Proves A \Sequent B$ असेल, तर
$\Lang L(C) \subseteq \Lang L(A) \cap \Lang L(B)$
या भाषा मधील असे सूत्र~$C$ असते की
$\Log{G3c} \Proves A \Sequent C$
आणि $\Log{G3c} \Proves C \Sequent B$.
\end{thm}

~\Log{G3c} मधील सिद्धता च्या रचनेचा
उपयोग करण्यासाठी आपण अंतर्वेशन प्रमेयाचे
अधिक व्यापक विधान सिद्ध करू.
त्यासाठी क्रमवर्तींचे पूर्वांग व उत्तरांग
प्रत्येकी दोन भागांत विभागलेले मानू.
$\maeh{\Gamma_1}{\Gamma_2} \Sequent
\maeh{\Delta_1}{\Delta_2}$ या लेखनात
$\Gamma_1$ आणि $\Delta_1$ पहिल्या भागात,
तर $\Gamma_2$ आणि $\Delta_2$ दुसऱ्या
भागात आहेत. पुढील उपप्रमेयावरून
अंतर्वेशन प्रमेय लगेच मिळते.

\begin{lem}[माएहाराचे उपप्रमेय]\label{pt:cut:itp:lem:maehara}
$\Log{G3c}
\Proves \maeh{\Gamma_1}{\Gamma_2} \Sequent \maeh{\Delta_1}{\Delta_2}$
असेल, तर $\Lang L(C)
\subseteq \Lang L(\Gamma_1,
\Delta_1) \cap \Lang L(\Gamma_2, \Delta_2)$
या भाषा मधील असे सूत्र~$C$ असते की
$\Log{G3c}
\Proves \Gamma_1 \Sequent \Delta_1, C$ आणि
$\Log{G3c} \Proves C,
\Gamma_2 \Sequent \Delta_2$.
\end{lem}
चारही संदर्भांतील सूत्रे वाक्ये असतील, तर अंतर्वेशकातील
मुक्त चरे सार्वत्रिकपणे बांधून वाक्यरूप अंतर्वेशक घेता येतो.
एका मुक्त चरासाठी पहिल्या सिद्धतेत तो चर आधी नव्या
स्थिरांकाने बदलून $\RightR{\lforall}$ लावायचा; बंद संदर्भांमुळे
आयगेन चराची अट पूर्ण होते. दुसऱ्या सिद्धतेत त्या चराचे
पद वापरून $\LeftR{\lforall}$ आणि आवश्यक दुर्बलीकरण लावायचे.
प्रत्येक मुक्त चरासाठी हे केल्याने दोन्ही क्रमवर्ती सिद्धतायोग्य
राहतात आणि स्थिरांक-विधेयांची भाषा-सीमा बदलत नाही.

\begin{proof}
  $\pi \Proves \maeh{\Gamma_1}{\Gamma_2} \Sequent
  \maeh{\Delta_1}{\Delta_2}$ असे गृहीत धरा.
  $n=\pheight{\pi}$ यावर विगमन करून विधान सिद्ध करू.

  $n=0$ असेल, तर
  $\maeh{\Gamma_1}{\Gamma_2} \Sequent
  \maeh{\Delta_1}{\Delta_2}$ हे स्वयंसिद्धक आहे
  आणि एखादे अणुरूप सूत्र~$A$
  हे $\Gamma_1, \Gamma_2$ आणि
  $\Delta_1, \Delta_2$ या दोन्ही बाजूंना येते.
  डावीकडे आणि उजवीकडे $A$ पहिल्या की
  दुसऱ्या भागात आहे यावर चार प्रसंग ठरतात.
  \begin{enumerate}
    \item $A \in \Gamma_1$ आणि $A \in \Delta_1$.
    $\Gamma_1 \Sequent \Delta_1, C$ आणि
    $C, \Gamma_2 \Sequent \Delta_2$ या दोन्ही
    सिद्धतायोग्य असतील असे~$C$ हवे.
    $A$ दोन्ही बाजूंना असल्याने, $C$ काहीही
    निवडले तरी पहिली क्रमवर्ती आधीच स्वयंसिद्धक आहे.
    $C, \Gamma_2 \Sequent \Delta_2$ ही दुसरी
    क्रमवर्ती सिद्धतायोग्य असण्याची खात्री
    करण्यासाठी $C \ident \lfalse$ घ्यावे लागते.
    \item $A \in \Gamma_1$ आणि $A \in \Delta_2$.
    मग $\Gamma_1 \Sequent \Delta_1, A$ आणि
    $A, \Gamma_2 \Sequent \Delta_2$ ही दोन्ही
    स्वयंसिद्धके आहेत; म्हणून $C \ident A$ घेता येते.
    \item $A \in \Gamma_2$ आणि $A \in \Delta_1$.
    मग $A, \Gamma_1 \Sequent \Delta_1$ आणि
    $\Gamma_2 \Sequent \Delta_2, A$ ही स्वयंसिद्धके
    असतील; पण त्यांत~$A$ चुकीच्या बाजूंना आहे.
    म्हणून $A$ अंतर्वेशक नाही.
    तथापि $\RightR{\lnot}$ आणि $\LeftR{\lnot}$
    वापरून $\Gamma_1
    \Sequent \Delta_1, \lnot A$ आणि
    $\lnot A, \Gamma_2 \Sequent
    \Delta_2$ या क्रमवर्ती सिद्ध करता येतात.
    \item $A \in \Gamma_2$ आणि $A \in \Delta_2$.
    हा सरावासाठीचा प्रसंग आहे.
  \end{enumerate}
  $\lfalse \in \Gamma_1$ किंवा
  $\lfalse \in \Gamma_2$ यांमुळे
  $\maeh{\Gamma_1}{\Gamma_2} \Sequent
  \maeh{\Delta_1}{\Delta_2}$ स्वयंसिद्धक असण्याचे
  प्रसंगही सरावासाठी सोडले आहेत.

  $n>0$ असेल, तर $\pi$ च्या शेवटी
  तार्किक अनुमान आहे. कोणता नियम वापरला
  आणि मुख्य सूत्र कोणत्या भागात आहे, यानुसार
  प्रसंग वेगळे करावे लागतील. प्रत्येक प्रसंगात
  $\pheight{\pi_1} < n$ असलेल्या सर्व
  सिद्धता~$\pi_1$ साठी, त्यांची अंतिम
  क्रमवर्ती कोणत्याही प्रकारे दोन भागांत
  विभागली असली तरी, उपप्रमेय खरे आहे हे
  विगमन गृहीतक वापरता येते. काही विशिष्ट
  प्रसंग पाहू.
  \begin{enumerate}
    \item $\pi$ च्या शेवटी $\RightR{\lnot}$
    असून मुख्य सूत्र~$\lnot
    A$ आहे. $\lnot A$ पहिल्या की
    दुसऱ्या भागात आहे, यानुसार $\pi$ चा
    शेवट पुढीलपैकी एक असतो:
    \[
    \AxiomC{}
    \RightLabel{$\pi_1$}
    \Deduce$\maeh{A, \Gamma_1}{\Gamma_2} \fCenter
      \maeh{\Delta_1}{\Delta_2}$
    \RightLabel{\RightR{\lnot}}
    \UnaryInf$\maeh{\Gamma_1}{\Gamma_2} \fCenter
      \maeh{\Delta_1, \lnot A}{\Delta_2}$
    \DisplayProof
    \]
    किंवा
    \[
    \AxiomC{}
    \RightLabel{$\pi_1$}
    \Deduce$\maeh{\Gamma_1}{A, \Gamma_2} \fCenter
      \maeh{\Delta_1}{\Delta_2}$
    \RightLabel{\RightR{\lnot}}
    \UnaryInf$\maeh{\Gamma_1}{\Gamma_2} \fCenter
      \maeh{\Delta_1}{\Delta_2, \lnot A}$
    \DisplayProof
    \]
    मुख्य सूत्र~$\lnot A$ पहिल्या
    भागात असेल, तर आधारक्रमवर्तीतील सक्रिय
    सूत्र~$A$ देखील आपण पहिल्या
    भागातच ठेवले आहे. ही निवड आपल्याला
    करता येते. सक्रिय सूत्र दुसऱ्या भागात
    ठेवले तरी $\pi_1$ च्या
    अंतिम क्रमवर्तीला विगमन गृहीतक लागू होईल;
    मात्र मग अंतर्वेशक शोधता येणार नाही
    (हे करून पाहा).

    आधी $\lnot A$ पहिल्या भागात
    असलेला प्रसंग पाहू. विगमन गृहीतकानुसार
    $\pi_1$ च्या अंतिम क्रमवर्तीला
    अंतर्वेशक असतो, म्हणजे असे
    सूत्र~$C_1$ असते की पुढील दोन्ही
    क्रमवर्ती
    \begin{align*}
    A, \Gamma_1 & \Sequent \Delta_1, C_1 &&\text{आणि} &
    C_1, \Gamma_2 & \Sequent \Delta_2
    \intertext{सिद्ध करता येतात. आता आपल्याला असे सूत्र~$C$ हवे की
    पुढील दोन्ही सिद्धतायोग्य असतील:}
    \Gamma_1 & \Sequent \Delta_1, \lnot A, C &&\text{आणि} &
    C, \Gamma_2 & \Sequent \Delta_2
    \end{align*}
    स्पष्टपणे $C \ident C_1$ घेता येते:
    उजवीकडील क्रमवर्ती सिद्धतायोग्य आहे
    हे विगमन गृहीतकाने आधीच मिळते, आणि
    डावीकडील क्रमवर्ती विगमन गृहीतकातून
    मिळालेल्या क्रमवर्तीवर~$\RightR{\lnot}$
    लावून मिळते. विगमन गृहीतकाने
    $\Lang L(C_1) \subseteq \Lang L(A, \Gamma_1,
    \Delta_1) \cap \Lang L(\Gamma_2, \Delta_2)$
    हेदेखील मिळते. $\Lang L(\lnot
    A) = \Lang L(A)$ आणि
    $\Lang L(C) = \Lang L(C_1)$ असल्याने
    $\Lang L(C_1) \subseteq \Lang L(\Gamma_1, \Delta_1, \lnot A) \cap \Lang
    L(\Gamma_2, \Delta_2)$.

    दुसऱ्या प्रसंगात $\lnot A$
    दुसऱ्या भागात असल्याने स्थिती थोडी
    वेगळी आहे. आधारक्रमवर्तीतील सक्रिय
    सूत्र देखील दुसऱ्या भागात
    ठेवून विगमन गृहीतक लावतो; त्यातून
    \begin{align*}
    \Gamma_1 & \Sequent \Delta_1, C_1 &&\text{आणि} &
    C_1, A, \Gamma_2 & \Sequent \Delta_2
    \intertext{सिद्ध करता येतात. यांतील दुसऱ्या क्रमवर्तीवर पुन्हा \RightR{\lnot} लावल्यास पुढील सिद्धता मिळतात:}
    \Gamma_1 & \Sequent \Delta_1, C_1 &&\text{आणि} &
    C_1, \Gamma_2 & \Sequent \Delta_2, \lnot A
    \end{align*}
    $C_1$ अंतर्वेशक असेल तर नेमक्या
    याच क्रमवर्ती सिद्धतायोग्य हव्या होत्या;
    म्हणून पुन्हा $C \ident C_1$ घेऊ.
    पूर्वीप्रमाणे $\Lang
    L(C_1) \subseteq \Lang L(\Gamma_1, \Delta_1) \cap \Lang L(\Gamma_2,
    \Delta_2, \lnot A)$.
    \item $\pi$ च्या शेवटी $\RightR{\lif}$
    असून मुख्य सूत्र~$A
    \lif B$ आहे. $A \lif
    B$ पहिल्या की दुसऱ्या भागात आहे
    यावर पुन्हा दोन प्रसंग मिळतात.
    सिद्धता~$\pi$ चा शेवट असा असतो:
    \[
    \AxiomC{}
    \RightLabel{$\pi_1$}
    \Deduce$\maeh{A, \Gamma_1}{\Gamma_2} \fCenter
      \maeh{\Delta_1, B}{\Delta_2}$
    \RightLabel{\RightR{\lif}}
    \UnaryInf$\maeh{\Gamma_1}{\Gamma_2} \fCenter
      \maeh{\Delta_1, A \lif B}{\Delta_2}$
    \DisplayProof
    \]
    किंवा
    \[
    \AxiomC{}
    \RightLabel{$\pi_1$}
    \Deduce$\maeh{\Gamma_1}{A, \Gamma_2} \fCenter
      \maeh{\Delta_1}{\Delta_2, B}$
    \RightLabel{\RightR{\lif}}
    \UnaryInf$\maeh{\Gamma_1}{\Gamma_2} \fCenter
      \maeh{\Delta_1}{\Delta_2, A \lif B}$
    \DisplayProof
    \]
    येथेही आधारक्रमवर्तीतील सक्रिय
    सूत्रे निष्कर्षातील मुख्य
    सूत्राच्याच भागात ठेवली आहेत.
    पहिल्या प्रसंगात विगमन गृहीतकाने
    पुढील सिद्धता मिळतात:
    \begin{align*}
    A, \Gamma_1 & \Sequent \Delta_1, B, C_1 &&\text{आणि} & C_1, \Gamma_2 & \Sequent \Delta_2
    \intertext{डावीकडील क्रमवर्ती \RightR{\lif} साठी योग्य आधारक्रमवर्ती आहे. म्हणून पुढील क्रमवर्ती सिद्धतायोग्य आहेत:}
    \Gamma_1 & \Sequent \Delta_1, A \lif B, C_1 &&\text{आणि} & C_1, \Gamma_2 & \Sequent \Delta_2
    \end{align*}
    म्हणून $C_1$ हा $\pi$ च्या
    अंतिम क्रमवर्तीचाही अंतर्वेशक आहे;
    पुन्हा $C \ident C_1$ घेता येते.

    दुसरा प्रसंगही याच प्रकारे हाताळता येतो.
    \item $\pi$ च्या शेवटी $\LeftR{\lif}$
    असून मुख्य सूत्र~$A
    \lif B$ आहे. $A \lif B$
    पहिल्या भागात असेल, तर
    सिद्धता~$\pi$ चा शेवट असा असतो:
    \[
    \AxiomC{}
    \RightLabel{$\pi_1$}
    \Deduce$\maeh{\Gamma_1}{\Gamma_2} \fCenter
      \maeh{\Delta_1, A}{\Delta_2}$
    \AxiomC{}
    \RightLabel{$\pi_2$}
    \Deduce$\maeh{B, \Gamma_1}{\Gamma_2} \fCenter
      \maeh{\Delta_1}{\Delta_2}$
    \BinaryInf$\maeh{A \lif B, \Gamma_1}{\Gamma_2} \fCenter
      \maeh{\Delta_1}{\Delta_2}$
    \DisplayProof
    \]
    विगमन गृहीतकाने आता चार
    सिद्धतायोग्य क्रमवर्ती मिळतात:
    डाव्या आधारक्रमवर्तीच्या
    अंतर्वेशक~$C_1$ साठी दोन आणि
    दुसऱ्या आधारक्रमवर्तीच्या
    अंतर्वेशक~$C_2$ साठी दोन:
    \begin{align*}
    \Gamma_1 & \Sequent \Delta_1, A, C_1 &&\text{(1)} &
    C_1, \Gamma_2 & \Sequent \Delta_2 && \text{(2)}\\
    B, \Gamma_1 & \Sequent \Delta_1, C_2 &&\text{(3)} &
    C_2, \Gamma_2 & \Sequent \Delta_2 && \text{(4)}
    \end{align*}
    दुर्बलीकरण लावून (1) च्या
    उजवीकडे~$C_2$ आणि (3) च्या
    उजवीकडे~$C_1$ जोडल्यास
    (1) आणि (3) या क्रमवर्ती
    \LeftR{\lif} च्या आधारक्रमवर्ती
    म्हणून वापरता येतात:
    \[
    \AxiomC{}
    \Deduce$\Gamma_1 \fCenter \Delta_1, A, C_1, C_2$
    \AxiomC{}
    \Deduce$B, \Gamma_1 \fCenter \Delta_1, C_1, C_2$
    \RightLabel{\LeftR{\lif}}
    \BinaryInf$A \lif B, \Gamma_1 \fCenter \Delta_1, C_1, C_2$
    \DisplayProof
    \]
    $\RightR{\lor}$ लावल्याने पुढील
    क्रमवर्ती मिळते:
    \[
    A \lif B, \Gamma_1 \fCenter \Delta_1, C_1 \lor C_2
    \]
    त्यामुळे या प्रसंगात
    $C_1 \lor C_2$ अंतर्वेशक असू
    शकतो असे सुचते. उरलेल्या
    (2) आणि~(4) क्रमवर्ती वापरून
    आवश्यक दुसरी क्रमवर्तीही
    सिद्ध करता येते:
    \[
    \AxiomC{}
    \Deduce$C_1, \Gamma_2 \fCenter \Delta_2$
    \AxiomC{}
    \Deduce$C_2, \Gamma_2 \fCenter \Delta_2$
    \RightLabel{\LeftR{\lor}}
    \BinaryInf$C_1 \lor C_2, \Gamma_2 \fCenter \Delta_2$
    \DisplayProof
    \]
    $C_1 \lor C_2$ योग्य
    भाषेत आहे हे तपासणे उरते.
    विगमन गृहीतकानुसार
    \begin{align*}
    \Lang L(C_1) & \subseteq
    \Lang L(\Gamma_1, \Delta_1, A) \cap \Lang L(\Gamma_2, \Delta_2) &\text{आणि}\\
    \Lang L(C_2) & \subseteq
    \Lang L(B, \Gamma_1, \Delta_1) \cap \Lang L(\Gamma_2, \Delta_2)
    \intertext{कारण}
    \Lang L(A) & \subseteq \Lang L(A \lif B) &\text{आणि}\\
    \Lang L(B) & \subseteq \Lang L(A \lif B)
    \intertext{म्हणून या दोन्ही अटी मिळतात}
    \Lang L(C_1) & \subseteq
    \Lang L(\Gamma_1, \Delta_1, A \lif B) \cap \Lang L(\Gamma_2, \Delta_2) &\text{आणि}\\
    \Lang L(C_2) & \subseteq
    \Lang L(\Gamma_1, \Delta_1, A \lif B) \cap \Lang L(\Gamma_2, \Delta_2)
    \intertext{आणि परिणामी, $\Lang L (C_1 \lor C_2) = {}$}
    \Lang L(C_1) \cup \Lang L(C_2) & \subseteq
    \Lang L(\Gamma_1, \Delta_1, A \lif B) \cap \Lang L(\Gamma_2, \Delta_2),
    \end{align*}
    आणि हेच अपेक्षित आहे.

    $A \lif B$ दुसऱ्या
    भागात असण्याचा प्रसंग
    वेगळा आहे, पण तशाच
    पद्धतीने हाताळता येतो.
    विगमन गृहीतकाने पुन्हा
    चार सिद्धतायोग्य क्रमवर्ती मिळतात:
    \begin{align*}
    \Gamma_1 & \Sequent \Delta_1, C_1 &&\text{(1)} &
    C_1, \Gamma_2 & \Sequent \Delta_2, A && \text{(2)}\\
    \Gamma_1 & \Sequent \Delta_1, C_2 &&\text{(3)} &
    C_2, B, \Gamma_2 & \Sequent \Delta_2 && \text{(4)}
    \end{align*}
    आता (2) आणि (4) या क्रमवर्ती,
    काही सूत्रे दुर्बल करून
    जोडल्यावर,~\LeftR{\lif}
    च्या आधारक्रमवर्ती ठरतात:
    \[
    \AxiomC{}
    \Deduce$C_1, C_2, \Gamma_2 \fCenter \Delta_2, A$
    \AxiomC{}
    \Deduce$B, C_1, C_2, \Gamma_2 \fCenter \Delta_2$
    \RightLabel{\LeftR{\lif}}
    \BinaryInf$A \lif B, C_1, C_2, \Gamma_2 \fCenter \Delta_2$
    \DisplayProof
    \]
    यांतून $C_1$ आणि $C_2$
    वापरून बनलेले एकच सूत्र
    असणारी क्रमवर्ती हवी आहे;
    पण यावेळी ते पूर्वांगात हवे
    (कारण आता दुसऱ्या भागाशी
    संबंध आहे). $\LeftR{\land}$
    लावल्याने ते साधते; म्हणून
    $C \ident
    (C_1 \land C_2)$ अंतर्वेशक घेऊ.
    उरलेल्या (1) आणि (3)
    क्रमवर्ती वापरून पहिल्या
    भागाची संबंधित क्रमवर्तीही
    सिद्ध करता येते:
    \[
    \AxiomC{}
    \Deduce$\Gamma_1 \fCenter \Delta_1, C_1$
    \AxiomC{}
    \Deduce$\Gamma_1 \fCenter \Delta_1, C_2$
    \RightLabel{\RightR{\land}}
    \BinaryInf$\Gamma_1 \fCenter \Delta_1, C_1 \land C_2$
    \DisplayProof
    \]
    पूर्वीप्रमाणे
    $\Lang L(C_1 \land C_2) \subseteq \Lang
    L(\Gamma_1, \Delta_1) \cap \Lang L(\Gamma_2, \Delta_2, A \lif
    B)$ मिळते.

    \item $\pi$ च्या शेवटी $\RightR{\lforall}$
    असून मुख्य सूत्र~$\lforall[x][A(x)]$ आहे:
    \[
    \AxiomC{}
    \RightLabel{$\pi_1$}
    \Deduce$\maeh{\Gamma_1}{\Gamma_2} \fCenter
    \maeh{\Delta_1, A(c)}{\Delta_2}$
    \RightLabel{\RightR{\lforall}}
    \UnaryInf$\maeh{\Gamma_1}{\Gamma_2} \fCenter
      \maeh{\Delta_1, \lforall[x][A(x)]}{\Delta_2}$
    \DisplayProof
    \]
    किंवा
    \[
    \AxiomC{}
    \RightLabel{$\pi_1$}
    \Deduce$\maeh{\Gamma_1}{\Gamma_2} \fCenter
      \maeh{\Delta_1}{\Delta_2, A(c)}$
    \RightLabel{\RightR{\lforall}}
    \UnaryInf$\maeh{\Gamma_1}{\Gamma_2} \fCenter
      \maeh{\Delta_1}{\Delta_2, \lforall[x][A(x)]}$
    \DisplayProof
    \]
    पहिल्या प्रसंगात विगमन गृहीतकाने
    $\Gamma_1 \Sequent \Delta_1, A(c), C_1$
    आणि $C_1, \Gamma_2
    \Sequent \Delta_2$ यांच्या सिद्धता मिळतात.
    पहिल्या क्रमवर्तीवर $\RightR{\lforall}$
    लावून $\Gamma_1 \Sequent \Delta_1, \lforall[x][A(x)], C_1$
    मिळते. ($\pi$ च्या अंतिम
    \RightR{\lforall} अनुमानात आयगेन चराची
    अट आधीच पूर्ण असल्याने येथेही ती पूर्ण आहे.)
    $\Lang
    L(C_1) \subseteq \Lang L(\Gamma_1, \Delta_1, \lforall[x][A(x)])
    \cap \Lang L(\Gamma_2, \Delta_2)$
    ही अट पूर्ण असेल, तर $C_1$ अंतर्वेशक ठरेल.
    विगमन गृहीतकाने
    $\Lang L(C_1) \subseteq \Lang L(\Gamma_1,
    \Delta_1, A(c)) \cap \Lang L(\Gamma_2, \Delta_2)$
    मिळते. म्हणून~$c$ ची काळजी घ्यावी लागते:
    $C_1$ मध्ये $c$ येऊ शकतो का? येऊ शकत नाही.
    $C_1$ मध्ये $c$ आल्यास
    $c \in \Lang L(\Gamma_1, \Delta_1,
    A(c))$ (जे खरे आहे) आणि
    $c \in \Lang L(\Gamma_2,
    \Delta_2)$ हेही दोन्ही खरे ठरतील.
    पण मग $\pi$ मधील
    $\RightR{\lforall}$ अनुमानाच्या
    निष्कर्षात $c$ येईल आणि
    आयगेन चराची अट मोडेल.

    दुसरा प्रसंगही यासारखाच आहे.

    \item $\pi$ च्या शेवटी
    $\RightR{\lexists}$ असून मुख्य
    सूत्र~$\lexists[x][A(x)]$ आहे.
    येथेही दोन प्रसंग आहेत.
    $\lexists[x][A(x)]$ पहिल्या
    भागात असलेला प्रसंग पाहू;
    दुसरा सरावासाठी सोडला आहे.

    पहिल्या प्रसंगात $\pi$ चा
    शेवट असा असतो:
    \[
    \AxiomC{}
    \RightLabel{$\pi_1$}
    \Deduce$\maeh{\Gamma_1}{\Gamma_2} \fCenter
      \maeh{\Delta_1, \lexists[x][A(x)], A(t)}{\Delta_2}$
    \RightLabel{\RightR{\lexists}}
    \UnaryInf$\maeh{\Gamma_1}{\Gamma_2} \fCenter
      \maeh{\Delta_1, \lexists[x][A(x)]}{\Delta_2}$
    \DisplayProof
    \]
    \begingroup\sloppy
    विगमन गृहीतकाने
    $\Gamma_1 \Sequent
    \Delta_1, \lexists[x][A(x)], A(t), C_1$
    आणि $C_1, \Gamma_2
    \Sequent \Delta_2$ यांच्या सिद्धता मिळतात.
    पहिल्या क्रमवर्तीवर $\RightR{\lexists}$
    लावल्याने $\Gamma_1 \Sequent \Delta_1, \lexists[x][A(x)], C_1$
    मिळते. म्हणून
    $\Lang L(C_1) \subseteq \Lang L(\Gamma_1, \Delta_1,
    \lexists[x][A(x)]) \cap \Lang L(\Gamma_2, \Delta_2)$
    ही अट पूर्ण असल्यास $C_1$ पुन्हा
    अंतर्वेशक ठरेल. येथे विगमन गृहीतकाने
    $\Lang L(C_1)
    \subseteq \Lang L(\Gamma_1, \Delta_1, \lexists[x][A(x)], A(t))
    \cap \Lang L(\Gamma_2, \Delta_2)$
    मिळते. आपण फलने नाहीत असे
    गृहीत धरले आहे, हे लक्षात घ्या.
    म्हणून पद~$t$ हे चर किंवा स्थिरांक असते.
    $t$ चर असेल, तर त्याच्या प्रतिस्थापनाने नवीन
    स्थिरांक किंवा विधेय-चिन्ह येत नाही:
    $\Lang L(A(t))\subseteq\Lang L(A(x))$.
    त्यामुळे वरील विगमन गृहीतकातील भाषा-सीमा
    अंतिम क्रमवर्तीच्या दोन भागांच्या भाषांच्या
    छेदातच राहते आणि $C\ident C_1$ घेता येते.
    $\RightR{\lexists}$ साठी आयगेन चराची अट नाही.
    आता $t$ हे स्थिरांक असल्याचा प्रसंग,
    म्हणजे $t\ident b$, पाहू.
    $b \notin \Lang L(C_1)$ असेल,
    किंवा $b \in \Lang L(\Gamma_1, \Delta_1, \lexists[x][A(x)]) \cap
    \Lang L(\Gamma_2, \Delta_2)$ असेल,
    तर अडचण नाही: $C_1$ अंतर्वेशक घेता येते.
    पण $b \in \Lang
    L(C_1)$ आणि $b \notin \Lang L(\Gamma_1, \Delta_1,
    \lexists[x][A(x)]) \cap \Lang L(\Gamma_2, \Delta_2)$
    असेल तर? प्रथम~$C_1$ मधील
    $b$ ची प्रत्येक उपस्थिती
    $C_1$ मध्ये न येणारा नवा चल~$y$
    निवडून त्याने बदलल्यावर
    $C_1 \ident \Subst{C_1'}{b}{y}$
    असे लिहिता येते; येथे $C_1'$
    मध्ये~$b$ येत नाही. शिवाय,
    $b \notin
    \Lang L(\Gamma_1, \Delta_1, \lexists[x][A(x)])$
    किंवा $b \notin
    \Lang L(\Gamma_2, \Delta_2)$.
    त्यानुसार $C \ident
    \lforall[y][C_1'(y)]$ किंवा
    $C \ident \lexists[y][C_1'(y)]$
    घेता येते. पहिल्या प्रसंगात:\par\endgroup
    \[
    \AxiomC{}
    \Deduce$\Gamma_1 \fCenter
      \Delta_1, \lexists[x][A(x)], A(b), C_1'(b)$
    \RightLabel{\RightR{\lexists}}
    \UnaryInf$\Gamma_1 \fCenter \Delta_1, \lexists[x][A(x)], C_1'(b)$
    \RightLabel{\RightR{\lforall}}
    \UnaryInf$\Gamma_1 \fCenter
      \Delta_1, \lexists[x][A(x)], \lforall[y][C_1'(y)]$
    \DisplayProof
    \]
    आणि
    \[
    \AxiomC{}
    \Deduce$C_1'(b), \lforall[y][C_1'(y)], \Gamma_2 \fCenter \Delta_2$
    \RightLabel{\LeftR{\lforall}}
    \UnaryInf$\lforall[y][C_1'(y)], \Gamma_2 \fCenter \Delta_2$
    \DisplayProof
    \]
    सर्वात वरच्या क्रमवर्तींच्या
    सिद्धता विगमन गृहीतकातून मिळतात
    (दुसऱ्या सिद्धतेत पूर्वांगात
    $\lforall[y][C_1'(y)]$ जोडण्यासाठी
    दुर्बलीकरणाची ग्राह्यता वापरतो).
    पहिल्या सिद्धता मधील
    \RightR{\lforall} साठी आयगेन चराची
    अट पूर्ण आहे, कारण
    $b \notin \Lang L(\Gamma_1, \Delta_1,
    \lexists[x][A(x)])$ आणि
    $C_1'$ मध्ये~$b$ येत नाही,
    अशीच त्याची व्याख्या केली आहे.

    दुसऱ्या प्रसंगात, विगमन गृहीतक
    आणि $\lexists[y][C_1'(y)]$
    जोडणाऱ्या दुर्बलीकरणाच्या
    ग्राह्यतेने पुढील मिळते:
    \[
    \AxiomC{}
    \Deduce$\Gamma_1 \fCenter
      \Delta_1, \lexists[x][A(x)], A(b), \lexists[y][C_1'(y)], C_1'(b)$
    \RightLabel{\RightR{\lexists}}
    \UnaryInf$\Gamma_1 \fCenter
      \Delta_1, \lexists[x][A(x)], \lexists[y][C_1'(y)], C_1'(b)$
    \RightLabel{\RightR{\lexists}}
    \UnaryInf$\Gamma_1 \fCenter
      \Delta_1, \lexists[x][A(x)], \lexists[y][C_1'(y)]$
    \DisplayProof
    \]
    आणि
    \[
    \AxiomC{}
    \Deduce$C_1'(b), \Gamma_2 \fCenter \Delta_2$
    \RightLabel{\LeftR{\lexists}}
    \UnaryInf$\lexists[y][C_1'(y)], \Gamma_2 \fCenter \Delta_2$
    \DisplayProof
    \]
    दुसऱ्या सिद्धता मधील
    \LeftR{\lexists} साठी आयगेन चराची
    अट पूर्ण आहे, कारण
    $b \notin \Lang L(\Gamma_2,
    \Delta_2)$ आणि $C_1'$ मध्ये~$b$
    येत नाही, अशीच त्याची व्याख्या केली आहे.
  \end{enumerate}
  उरलेले प्रसंगही यांसारखे आहेत;
  ते सरावासाठी सोडले आहेत.
\end{proof}

\begin{prob}
$\lnand$ हा संयोजक पुढील नियमांसह
आपल्याकडे आहे असे समजा:
\[
\Axiom$\Gamma \fCenter \Delta, A$
\Axiom$\Gamma \fCenter \Delta, B$
\RightLabel{\LeftR{\lnand}}
\BinaryInf$A \lnand B, \Gamma \fCenter \Delta$
\DisplayProof
\quad
\Axiom$A, B, \Gamma \fCenter \Delta$
\RightLabel{\RightR{\lnand}}
\UnaryInf$\Gamma \fCenter \Delta, A \lnand B$
\DisplayProof
\]
ज्या प्रसंगात $\pi$ चा शेवट
यांपैकी एखाद्या नियमाने होतो,
ते \ref{pt:cut:itp:lem:maehara} च्या सिद्धतेत
पूर्ण करून माएहाराचे उपप्रमेय
अजूनही खरे आहे हे दाखवा.
\end{prob}

माएहाराचे उपप्रमेय सिद्ध झाल्यावर
त्यावरून क्रेगचे अंतर्वेशन प्रमेय
सिद्ध करता येते.

\begin{proof}[\ref{pt:cut:itp:thm:interpolation} ची सिद्धता]
  $A \Sequent B$ सिद्धतायोग्य आहे असे
  समजा. $\maeh{A}{\ } \Sequent \maeh{\ }{B}$
  याला \ref{pt:cut:itp:lem:maehara} लावून
  अंतर्वेशक~$C$ मिळवा. मग
  $\Proves A \Sequent C$ आणि
  $\Proves C
  \Sequent B$.
\end{proof}

$\Gamma$ ही वाक्यांचा संच असलेली
उपपत्ती, भाषा~$\Lang L$ मध्ये आहे
आणि त्यात $n$-स्थानी {विधेय}~$R$ आहे
असे समजा. पुढील अट पूर्ण असेल, तेव्हा
आणि केवळ तेव्हाच, $\Gamma$ ही
$R$ ला \emph{अंतर्निहितपणे परिभाषित करते}:
\begin{align*}
\Gamma, \Gamma'
& \Entails \lforall[x_1][\dots\lforall[x_n][(R(x_1,\dots,x_n) \liff
R'(x_1, \dots, x_n))]],
\intertext{येथे $\Gamma'$ म्हणजे $\Gamma$ मध्ये $R$ ची प्रत्येक उपस्थिती
$R' \notin \Lang L$ ने बदलल्यावर मिळणारी उपपत्ती. $\Gamma$ ही $R$ ला
\emph{स्पष्टपणे परिभाषित करते}, तेव्हा आणि केवळ तेव्हाच, $R$ नसलेले असे
$A(x_1, \dots, x_n)$ असते की}
\Gamma &\Entails
\lforall[x_1][\dots\lforall[x_n][(R(x_1,\dots,x_n) \liff A(x_1,
\dots, x_n))]].
\end{align*}
स्पष्टपणे, $\Gamma$ ही $R$ ला
स्पष्टपणे परिभाषित करत असेल, तर
ती $R$ ला अंतर्निहितपणेही परिभाषित करते.
उलट विधान लगेच स्पष्ट होत नसले,
तरी ते खरे आहे:

\begin{lem}[बेथचे परिभाष्यता प्रमेय]
  $\Gamma$ ही $R$ ला अंतर्निहितपणे
  परिभाषित करत असेल, तर ती त्याला
  स्पष्टपणेही परिभाषित करते.
\end{lem}

\begin{proof}
  $\Gamma$ ही $R$ ला अंतर्निहितपणे
  परिभाषित करते असे समजा. वरीलप्रमाणे
  $R'$ आणि $\Gamma'$ घ्या;
  $\Lang L' = \Lang L(\Gamma') = \Lang L \setminus \{R\}
  \cup \{R'\}$ असू द्या. गृहीतकानुसार
  \begin{align*}
    \Gamma, \Gamma' &
    \Sequent \lforall[x_1][\dots\lforall[x_n][(R(x_1,\dots,x_n)
    \liff R'(x_1, \dots, x_n))]]
  \intertext{$c_1$, \dots,~$c_n$ हे $\Gamma$ मध्ये न येणारे स्थिरांक असू द्या. व्युत्क्रमणीयतेच्या उपप्रमेयाने मिळते:}
    \Gamma, \Gamma', R(c_1,\dots,c_n) & \Sequent  R'(c_1, \dots, c_n)
  \intertext{माएहाराचे उपप्रमेय पुढील क्रमवर्तीला लावा:}
    \maeh{\Gamma, R(c_1,\dots,c_n)}{\Gamma'} & \Sequent
      \maeh{\ }{R'(c_1, \dots, c_n)}
  \intertext{या बंद संदर्भांसाठी आवश्यक मुक्त चरे सार्वत्रिकपणे बांधून वाक्यरूप अंतर्वेशक निवडा. मग नव्या स्थिरांकांच्या जागी संबंधित चरे ठेवून असे~$A \ident A(x_1, \dots, x_n)$ मिळते की}
    \Gamma, R(c_1,\dots,c_n) & \Sequent A(c_1, \dots, c_n) \\
    A(c_1, \dots, c_n), \Gamma' & \Sequent R'(c_1, \dots, c_n)
  \intertext{यांच्या सिद्धता असतात. $\Lang L(A) \subseteq
  (\Lang L(\Gamma) \setminus \{R\}) \cup \{c_1,\dots,c_n\}$ असल्याने $A$ मध्ये~$R$ नसतो.
  दुसऱ्या क्रमवर्तीच्या सिद्धता मध्ये सर्वत्र $R'$ च्या जागी~$R$ ठेवून
  पुढील सिद्धता मिळते:}
    A(c_1, \dots, c_n), \Gamma & \Sequent R(c_1, \dots, c_n)
  \end{align*}
  येथील नव्या स्थिरांकांच्या जागी
  संबंधित चरे ठेवता येतात.
  \RightR{\lif} आणि
  \RightR{\lforall} लावून पुढील
  सिद्धता मिळते:
  \[
    \Gamma \Sequent \lforall[x_1][\dots\lforall[x_n][(R(x_1,\dots,x_n) \liff A(x_1, \dots, x_n))]].
  \]
  म्हणून $A(x_1, \dots, x_n)$ हे
  $\Gamma$ मध्ये $R$ ला स्पष्टपणे
  परिभाषित करते.
\end{proof}

अंतर्वेशनाशी समतुल्य असलेला
आणि त्याप्रमाणेच माएहाराचे उपप्रमेय वापरून
सिद्ध करता येणारा तिसरा निष्कर्ष असा:
दोन सुसंगत उपपत्ती~$\Gamma_1$ आणि
$\Gamma_2$ यांत $\Lang L(\Gamma) = \Lang L(\Gamma_1)
\cap \Lang L(\Gamma_2)$ या भाषेतील
पूर्ण उपपत्ती~$\Gamma$ समाविष्ट असेल,
तर $\Gamma_1 \cup \Gamma_2$ सुसंगत असते.

पूर्ण उपपत्ती म्हणजे तिच्या भाषेतील
प्रत्येक वाक्य~$A$ साठी
$\Gamma \Proves A$ किंवा
$\Gamma \Proves \lnot A$ सिद्ध होणारी उपपत्ती,
हे आठवा. $\Gamma_1$ आणि~$\Gamma_2$
या $\Gamma$ च्या सुसंगत विस्तारे
असल्याने $\Gamma$ सुसंगतच असली पाहिजे.
म्हणून $A$ हे
$\Lang L(\Gamma_1) \cap \Lang
L(\Gamma_2)$ या भाषेतील
वाक्य असेल, तर
$\Gamma_1 \Entails A$ आणि
$\Gamma_2
\Entails \lnot A$ दोन्ही शक्य नाहीत.
तसे एखादे~$A$ आहे असे समजा.
$\Gamma_1 \Entails A$ असेल, तर
$\Gamma \Entails A$.
अन्यथा, $\Gamma$ पूर्ण असल्याने
$\Gamma \Entails \lnot A$;
आणि $\Gamma \subseteq \Gamma_1$
असल्याने $\Gamma_1 \Entails \lnot A$,
म्हणजे $\Gamma_1$ विसंगत ठरेल.
म्हणून $\Gamma \Entails A$;
आता $\Gamma
\subseteq \Gamma_2$ असल्याने
$\Gamma_2 \Entails A$.
त्यामुळे $\Gamma \Entails/
\lnot A$; अन्यथा $\Gamma_2$
विसंगत ठरेल.

संयुक्त भाषा~
$\Lang L(\Gamma_1) \cap \Lang L(\Gamma_2)$
मधील असे वाक्य~$A$
नसणे की $\Gamma_1
\Entails A$ आणि
$\Gamma_2 \Entails \lnot A$,
हेच सिद्धतेसाठी पुरेसे आहे.
आता माएहाराचे उपप्रमेय वापरून
ही सिद्धता सिद्धता-उपपत्तीय
पद्धतीने देऊ.

\begin{thm}[रॉबिन्सनचे संयुक्त सुसंगतता प्रमेय]
$\Gamma_1$ आणि $\Gamma_2$
या सुसंगत उपपत्ती असू द्या.
$\Lang L(\Gamma_1) \cap \Lang L(\Gamma_2)$
या भाषा मधील कोणत्याही~$A$
साठी $\Gamma_1 \Entails A$
आणि $\Gamma_2 \Entails \lnot A$
दोन्ही होत नसतील, तर
$\Gamma_1 \cup \Gamma_2$ सुसंगत असते.
\end{thm}

\begin{proof}
  उलट, $\Gamma_1 \cup \Gamma_2$
  विसंगत आहे असे समजा.
  मग $\Gamma_1, \Gamma_2 \Sequent \ $
  या क्रमवर्तीची सिद्धता आहे.
  $\maeh{\Gamma_1}{\Gamma_2}
  \Sequent \maeh{\ }{\ }$ याला
  माएहाराचे उपप्रमेय लावून, आवश्यक असल्यास वर
  सांगितल्याप्रमाणे अंतर्वेशकातील मुक्त चरे बांधल्याने,
  $\Lang L(\Gamma_1) \cap \Lang L(\Gamma_2)$
  या भाषा मधील असे
  वाक्य~$A$ मिळते की
  $\Proves \Gamma_1 \Sequent A$
  आणि $\Proves A, \Gamma_2 \Sequent
  \quad$. दुसऱ्या क्रमवर्तीवरून
  $\Proves \Gamma_2 \Sequent \lnot A$
  मिळते. पण असे वाक्य
  अस्तित्वात नाही, हेच आपण
  गृहीत धरले होते.
\end{proof}

वर आपण उपपत्ती $\Gamma$,
$\Gamma_1$ आणि $\Gamma_2$
सांत आहेत असे नकळत गृहीत धरले.
संहततेमुळे हे निष्कर्ष
अनंत उपपत्तींनाही लागू होतात.



